\documentclass[11pt,reqno]{amsart}

\usepackage[utf8]{inputenc}
\usepackage[T1]{fontenc}
\usepackage{amsmath,amssymb,amsthm,mathtools}
\usepackage{array}
\usepackage[margin=1in]{geometry}
\usepackage{xcolor}
\usepackage[colorlinks=true,linkcolor=blue!50!black,citecolor=blue!50!black]{hyperref}

\theoremstyle{plain}
\newtheorem{theorem}{Theorem}[section]
\newtheorem{proposition}[theorem]{Proposition}
\newtheorem{lemma}[theorem]{Lemma}
\newtheorem{corollary}[theorem]{Corollary}

\theoremstyle{definition}
\newtheorem{definition}[theorem]{Definition}
\newtheorem{remark}[theorem]{Remark}

\DeclareMathOperator{\Haar}{Haar}
\newcommand{\chiad}{\chi_{\mathrm{ad}}}
\newcommand{\Qcone}{\mathcal Q}
\newcolumntype{L}[1]{>{\raggedright\arraybackslash}p{#1}}

\title[Full adjoint-generated Ginibre $Q_3$]{The full adjoint-generated case of Ginibre's condition \texorpdfstring{$Q_3$}{Q3}\\
\large Part III: completion of the full signed-product hierarchy}
\author{Leslie P. Polzer}
\address{}
\email{polzer@fastmail.com}
\date{July 19, 2026}

\hypersetup{
  pdftitle={The full adjoint-generated case of Ginibre's condition Q3},
  pdfauthor={Leslie P. Polzer},
  pdfsubject={Full signed-product hierarchy on the adjoint-generated cone},
  pdfkeywords={Ginibre Q3, compact Lie group, adjoint character, exact certificates}
}

\begin{document}

\begin{abstract}
Let $G$ be a compact connected Lie group with simple Lie algebra and let
$\Qcone_G$ be the closed multiplicative convex cone generated by its real
adjoint character.  We reduce Ginibre's full condition~$Q_3$ on $\Qcone_G$
to one explicit two-parameter hierarchy and prove that hierarchy for every
compact connected $G$ with simple Lie algebra.  The proof treats every finite
tuple length, every sign assignment, and every even number of minus factors;
Ginibre's Proposition~1 then promotes the adjoint generator to $\Qcone_G$.
The last finite classical boxes are closed by bounded Littlewood determinant
identities and an exact prime-field/GMP certificate covering 12,993 residual
inequalities.  Here ``full'' refers to Ginibre's finite-tuple and sign-pattern
quantifiers on the explicitly stated adjoint-generated cone; it does not mean
the cone of all positive-definite functions.  We also give an exact
$\mathrm{SO}(3)$ counterexample showing
that Ginibre's abelian cone of all real positive-definite functions cannot be
carried over literally to nonabelian compact groups.
\end{abstract}

\maketitle

\begin{remark}[Submission architecture and audit status]
\label{rem:submission-architecture}
The formal proof has three numbered parts and one detailed supplement
collected in three files: \path{paper.pdf} consolidates Parts~I--II,
\path{paper_full.pdf} is their 503-page formal detailed supplement, and the
present file is Part~III.  All three files are load-bearing.  In particular,
the detailed supplement supplies the numbered $B/C$ correction-prefix proofs
used by the compact residual-closure proposition in \path{paper.pdf}.  The
concise dependency spine is recorded
in \path{PUBLICATION_PROOF_SPINE.md}; diagnostic and superseded routes in the
expanded archive are not inputs to that spine.  Files under \path{referee_audits/} and the
historical report in the source archive are author-controlled self-audits.
They are reproducibility aids, not independent peer review.  The theorem
depends only on the numbered results and authenticated computations mapped in
the code and data availability section.
\end{remark}

\section{The exact adjoint-generated target}

The cone studied here is the smallest uniformly closed cone containing the
adjoint character and closed under multiplication and nonnegative linear
combination.  It is natural for two independent reasons.  Representation
theoretically, its polynomial part consists exactly of nonnegative linear
combinations of tensor-power characters
\(1,\chiad,\chiad^2,\ldots\); hence its Haar moments are invariant-tensor
multiplicities in tensor powers of the adjoint module.  From the viewpoint of
Ginibre's Proposition~1, it is also the maximal cone to which a proof for the
single generator \(\chiad\) promotes without introducing a second generator
or a mixed-character inequality.  The theorem therefore isolates the full
signed-product consequences of the adjoint tensor category.  It does not
assert that a general central positive-definite function is a uniform limit
of these polynomials, and it supplies no conclusion for a model whose
interaction or observables do not belong to this cone.

Let $G$ be as in the abstract, let $\mu_G$ be normalized Haar measure, and
write $X=\chiad(g)$ and $Y=\chiad(h)$ for independent Haar variables.  Since
the adjoint representation is real, $X$ and $Y$ are real-valued.

\begin{definition}[Adjoint-generated cone]
The adjoint-generated multiplicative convex cone is
\[
 \Qcone_G=
 \overline{\left\{
   \sum_{j=0}^{d}c_j\chiad^j:
   d\geq0,\ c_j\geq0
 \right\}}^{\lVert\cdot\rVert_\infty}
 \subset C(G).
\]
For integers $a,n\geq0$, define
\[
 I^G_{a,n}
 =\int_{G\times G}
   (\chiad(g)-\chiad(h))^{2a}
   (\chiad(g)+\chiad(h))^n\,d\mu_G(g)d\mu_G(h).
\]
\end{definition}

\begin{lemma}[The adjoint-generated cone is positive definite]
\label{lem:adjoint-cone-positive-definite}
Every element of $\Qcone_G$ is a real continuous positive-definite function
on $G$.  In particular, $\Qcone_G$ is a subcone of the cone $\mathcal P_G$
used in Proposition~\ref{prop:all-pd-counterexample} below.
\end{lemma}

\begin{proof}
Choose an orthonormal basis $(e_j)$ of the real adjoint module.  Then
\[
 \chiad(g)=\sum_j\langle e_j,\operatorname{Ad}(g)e_j\rangle
\]
is positive definite: for arbitrary $g_1,\ldots,g_r\in G$ and
$z_1,\ldots,z_r\in\mathbb C$ its defining quadratic form is
\[
 \sum_j\left\lVert\sum_{k=1}^r
 z_k\operatorname{Ad}(g_k)e_j\right\rVert^2\geq0.
\]
Products of positive-definite functions are positive definite by tensoring
their unitary representations, and nonnegative sums are positive definite by
direct sums and scaling.  Finally, a uniform limit preserves every finite
positive-semidefinite Gram matrix entrywise.  The character and all its
powers are real and continuous, proving the assertion from the definition of
$\Qcone_G$.
\end{proof}

\begin{theorem}[Full $Q_3$ reduction on the adjoint-generated cone]
\label{thm:full-cone-reduction}
Ginibre's condition $Q_3$ holds for $\Qcone_G$ if and only if
\[
 I^G_{a,n}\geq0\qquad(a,n\geq0).
\]
The only nonautomatic members of this hierarchy are
\[
 a\geq1,\qquad n\ \text{odd}.
\]
Theorem~\ref{thm:two-minus-import} below is exactly the row $a=1$.
\end{theorem}

\begin{proof}
For the singleton generating set $S_G=\{\chiad\}$, a signed product with
$r$ minus factors and $n$ plus factors is
\[
 \int_{G\times G}(X-Y)^r(X+Y)^n\,d\mu_G\,d\mu_G.
\]
It vanishes when $r$ is odd, by interchanging $X$ and $Y$, and for $r=2a$
it is $I^G_{a,n}$.  Thus the displayed hierarchy is precisely condition
$Q_3$ for the generating set $S_G$.

For completeness, the passage from the generator to the cone uses the
factorization
\[
 \begin{aligned}
 (fg)(x)\pm(fg)(y)
 =\frac12\bigl[&(f(x)+f(y))(g(x)\pm g(y))\\
               &+(f(x)-f(y))(g(x)\mp g(y))\bigr].
 \end{aligned}
\]
Iterating this identity expands every signed factor belonging to a product
of generators as a sum, with nonnegative coefficients, of products of
signed generator factors.  Convex combinations preserve the inequality,
and uniform limits preserve it because Haar measure is finite.  This is the
argument of Ginibre~\cite[Proposition~1 and (1.5)]{Ginibre1970}; it proves
that $S_G$ satisfies $Q_3$ if and only if $\Qcone_G$ does.

It remains to identify the automatic sectors.  If $n$ is even, the
integrand defining $I^G_{a,n}$ is pointwise nonnegative.  When $a=0$, expand
\[
 \mathbf E(X+Y)^n
 =\sum_{j=0}^{n}\binom nj m_j^Gm_{n-j}^G,
 \qquad
 m_j^G=\dim\bigl((\mathfrak g_{\mathbb C}^{\otimes j})^G\bigr)\geq0.
\]
Hence the all-plus sector is nonnegative as well.
\end{proof}

\begin{theorem}[Two-minus adjoint-character theorem]
\label{thm:two-minus-import}
For every compact connected Lie group $G$ with simple Lie algebra and every
integer $n\geq0$,
\[
 I^G_{1,n}\geq0.
\]
\end{theorem}

\begin{proof}
This is the main theorem of the journal-facing Parts~I--II manuscript,
\cite[Theorem~2.2]{AdjointTwoMinus}.  In the notation of that theorem its
functional $Q_3^G(n)$ is exactly $I^G_{1,n}$ above.  That manuscript is
generated as \path{paper.pdf} from \path{paper.tex} and is submitted with the
present Part~III.  The formal detailed supplement built from the same source
supplies the numbered correction-prefix proofs used by the compact
manuscript.
\end{proof}

\begin{proposition}[The all-positive-definite analogue is false]
\label{prop:all-pd-counterexample}
Let $G=\mathrm{SO}(3)$ with normalized Haar measure, and let
$\mathcal P_G$ denote the cone of all real continuous positive-definite
functions on $G$.  Then $\mathcal P_G$ does not satisfy Ginibre's
condition $Q_3$.  More precisely, four functions in $\mathcal P_G$ and the
sign pattern $--++$ give the value
\[
 -\frac8{279}.
\]
Consequently the adjoint-generated restriction in
Theorem~\ref{thm:full-cone-reduction} cannot be removed, even for a compact
connected group with simple Lie algebra.
\end{proposition}

\begin{proof}
For a unit vector $v\in\mathbb R^3$, put
\[
 f_v(g)=\langle v,gv\rangle ,\qquad g\in\mathrm{SO}(3).
\]
This function is real and continuous.  It is positive definite because,
for arbitrary $g_1,\ldots,g_r\in G$ and $z_1,\ldots,z_r\in\mathbb C$,
\[
 \sum_{i,j=1}^r\overline{z_i}z_j f_v(g_i^{-1}g_j)
 =\left\lVert\sum_{j=1}^r z_jg_jv\right\rVert^2\geq0.
\]

We first record the Haar integrals needed below.  For unit vectors
$v_1,\ldots,v_4$, define
\[
 c_1=\langle v_1,v_2\rangle\langle v_3,v_4\rangle,
 \quad
 c_2=\langle v_1,v_3\rangle\langle v_2,v_4\rangle,
 \quad
 c_3=\langle v_1,v_4\rangle\langle v_2,v_3\rangle.
\]
Haar projection on the standard representation gives
\[
 \int_G f_{v_i}=0,
 \qquad
 \int_G f_{v_i}f_{v_j}
 =\frac{\langle v_i,v_j\rangle^2}{3}.
\]
For the fourth integral, the three pairing tensors in
$(\mathbb R^3)^{\otimes4}$ have Gram matrix with diagonal entries $9$ and
off-diagonal entries $3$.  Its inverse has diagonal entries $2/15$ and
off-diagonal entries $-1/30$.  Since Haar averaging is the orthogonal
projection onto the invariant tensors, contraction with
$v_1\otimes\cdots\otimes v_4$ on both sides gives
\begin{equation}
\label{eq:so3-four-coefficient-integral}
 \int_G\prod_{i=1}^4f_{v_i}
 =\frac2{15}\sum_{i=1}^3c_i^2
  -\frac1{15}\sum_{1\leq i<j\leq3}c_ic_j.
\end{equation}

Take
\[
 v_1=\left(\sqrt{\frac{19}{31}},\sqrt{\frac{12}{31}},0\right),
 \quad
 v_2=\left(\sqrt{\frac{19}{31}},-\sqrt{\frac{12}{31}},0\right),
 \quad
 v_3=v_4=(1,0,0).
\]
Then $(c_1,c_2,c_3)=(7,19,19)/31$, and
\eqref{eq:so3-four-coefficient-integral} equals $61/961$.  Expanding the
signed product and using the vanishing one-point integrals yields, with
$M_{ij}=\int_G f_{v_i}f_{v_j}\,d\mu_G$,
\[
\begin{aligned}
 \mathcal J&=\int_{G\times G}
 (f_{v_1}(g)-f_{v_1}(h))\\
 &\qquad{}\cdot(f_{v_2}(g)-f_{v_2}(h))\\
 &\qquad{}\cdot(f_{v_3}(g)+f_{v_3}(h))\\
 &\qquad{}\cdot(f_{v_4}(g)+f_{v_4}(h))\\
 &=2\left[\int_G\prod_{i=1}^4f_{v_i}\,d\mu_G
 +M_{12}M_{34}-M_{13}M_{24}-M_{14}M_{23}\right]\\
 &=2\left[\frac{61}{961}
  +\frac{49-361-361}{9\cdot961}\right]
 =-\frac8{279}.
\end{aligned}
\]
This is a strict violation of $Q_3$ inside $\mathcal P_G$.
\end{proof}

\begin{remark}[Comparison with Ginibre's abelian Example~4]
\label{rem:ginibre-example4-comparison}
Ginibre~\cite[Example~4]{Ginibre1970} proves $Q_3$ for the cone of all real
positive-definite functions when the compact group is a finite product of
circles and finite cyclic groups.  At the end of that paper he proposes the
same cone on a noncommutative compact group only as a candidate.  Proposition
\ref{prop:all-pd-counterexample} shows that this proposed extension is false.
The result proved here matches Ginibre's complete $Q_3$ quantifiers and his
Proposition~1 cone-promotion mechanism, but on the explicitly smaller
multiplicative cone generated by the adjoint character.
The matrix coefficients in Proposition~\ref{prop:all-pd-counterexample} are
not central.  Consequently neither that counterexample nor the affirmative
theorem decides $Q_3$ for the intermediate cone of all real continuous
central positive-definite functions on a nonabelian compact group.
\end{remark}

\begin{remark}[Relation to earlier nonabelian correlation inequalities]
\label{rem:related-nonabelian-inequalities}
Sylvester~\cite[Secs.~1, 3--4]{Sylvester1980} had already disproved the
general Ginibre inequality for higher-dimensional $O(N)$ spin systems by
constructing graph-cycle counterexamples.  His variables are spins on a
product of spheres and his signed factors are indexed by graph edges.  The
present Proposition~\ref{prop:all-pd-counterexample} supplies a different,
exact witness directly in the compact-group formulation proposed at the end
of Ginibre's 1970 paper: four explicit real positive-definite functions on
$\mathrm{SO}(3)$, integrated over two normalized Haar variables, give
$-8/279$.  We therefore make no claim that this is the first nonabelian
failure of a Ginibre inequality.

Abdesselam~\cite[Sec.~1 and Thms.~2.1--2.3]{Abdesselam2022} subsequently
formulated general and padded Ginibre inequalities for invariant observables
of $O(N)$ and $\mathbb{CP}^{N-1}$ models.  He proves large-power asymptotics
in terms of Kirchhoff polynomials and proves all padded Ginibre inequalities
for inverse half-integer powers of stable determinantal polynomials.  Those
results concern an asymptotic regime and a determinantal-polynomial
substitute; they do not establish the finite-rank compact-group Haar
hierarchy proved here.  Conversely, the present theorem is restricted to the
adjoint-generated cone and does not settle the invariant-observable GKS2
problem for finite $O(N)$ spin systems, which Abdesselam records as open for
$N\geq3$.
\end{remark}

\begin{remark}[Scope of potential quantum applications]
\label{rem:quantum-application-scope}
Theorem~\ref{thm:full-adjoint-generated-q3} is an algebraic Haar-positivity
theorem.  Through Ginibre's general mechanism it can serve as an input for
group-valued lattice models whose interactions genuinely lie in
$\Qcone_G$.  It does not by itself prove correlation inequalities for quantum
Heisenberg or general $O(N)$ vector-spin systems, nonlinear sigma models,
lattice gauge theories, or continuum quantum field theories.  In particular,
the occurrence of an adjoint representation in a gauge model neither places
its plaquette interactions and observables in $\Qcone_G$ nor removes the
shared-link constraints.  This distinction is visible in the recent work of
Chevyrev and Garban~\cite[Thm.~1.5, Cor.~1.6, and Rem.~1.7]{ChevyrevGarban2025}:
their Villain-action limit applies to nonabelian lattice gauge theories, while
the Ginibre monotonicity consequence established there is abelian.

Such applications require additional model-specific results: Gibbs-measure
promotion on the relevant observable cone, a thermodynamic-limit argument,
and, for a quantum reconstruction, a reflection-positive transfer-matrix or
equivalent construction.  A direct route toward the conventional finite-$N$
$O(N)$ problem would instead require the two-minus or padded inequality for
defining-representation matrix coefficients on products of spheres.  That is
a separate theorem and is not asserted here.
\end{remark}

\begin{lemma}[Moment and generating-function forms]\label{lem:full-moment-form}
For $a,n\geq0$,
\[
 I^G_{a,n}
 =\sum_{j=0}^{2a}\sum_{k=0}^{n}
   (-1)^j\binom{2a}{j}\binom nk
   m^G_{2a-j+k}m^G_{j+n-k}.
\]
If
\[
 M_G(t)=\mathbf E e^{tX}
       =\sum_{r\geq0}m_r^G\frac{t^r}{r!},
\]
then, as an identity of convergent power series,
\[
 F_a^G(t):=\mathbf E\bigl[(X-Y)^{2a}e^{t(X+Y)}\bigr]
 =\sum_{j=0}^{2a}(-1)^j\binom{2a}{j}
   M_G^{(2a-j)}(t)M_G^{(j)}(t)
 =\sum_{n\geq0}I^G_{a,n}\frac{t^n}{n!}.
\]
Thus the full hierarchy is equivalent to coefficientwise nonnegativity of
every $F_a^G$.
\end{lemma}

\begin{proof}
Expand $(X-Y)^{2a}(X+Y)^n$ twice by the binomial theorem and use the
independence of $X$ and $Y$.  Expanding only $(X-Y)^{2a}$ after inserting
$e^{tX}e^{tY}$ gives the derivative formula.  Finally, expansion of
$e^{t(X+Y)}$ gives the last equality.  The adjoint character is bounded on
the compact group, so all series converge absolutely and termwise
integration is justified.
\end{proof}

\begin{lemma}[Central-cover invariance]\label{lem:full-central-cover}
If $\pi:\widetilde G\to G$ is a finite central covering of compact connected
Lie groups, then
\[
 I^{\widetilde G}_{a,n}=I^G_{a,n}
 \qquad(a,n\geq0).
\]
Thus every result below depends only on the simple Lie algebra and applies to
all of its compact connected global forms.
\end{lemma}

\begin{proof}
The adjoint characters satisfy
$\chiad^{\widetilde G}=\chiad^G\circ\pi$, and the push-forward of normalized
Haar measure under $\pi$ is normalized Haar measure.  Apply these facts in
both variables of the defining integral.
\end{proof}

\begin{lemma}[Nonnegative cumulants give nonnegative moments]
\label{lem:nonnegative-cumulants}
Let
\[
 K(t)=\sum_{r\geq1}\kappa_r\frac{t^r}{r!}
 \quad\text{with}\quad \kappa_r\geq0,
\]
and write
\[
 e^{K(t)}=\sum_{n\geq0}q_n\frac{t^n}{n!}.
\]
Then $q_n\geq0$ for every $n$.
\end{lemma}

\begin{proof}
The exponential formula gives
\[
 q_n=\sum_{\substack{k_1,\ldots,k_n\geq0\\
                      k_1+2k_2+\cdots+nk_n=n}}
 n!\prod_{r=1}^{n}
 \frac1{k_r!}\left(\frac{\kappa_r}{r!}\right)^{k_r}.
\]
Every summand is nonnegative.
\end{proof}

\begin{lemma}[Adjoint-character normalization]\label{lem:adjoint-normalization}
For $G$ compact and connected with simple Lie algebra,
\[
 \int_G\chiad(g)\,d\mu_G(g)=0,
 \qquad
 \int_G\chiad(g)^2\,d\mu_G(g)=1,
 \qquad
 \max_G\chiad=\dim\mathfrak g.
\]
\end{lemma}

\begin{proof}
The complexified adjoint representation is irreducible: a
$G$-invariant complex subspace differentiates, by connectedness, to an
$\operatorname{ad}(\mathfrak g_{\mathbb C})$-invariant subspace and hence is
an ideal of the simple Lie algebra $\mathfrak g_{\mathbb C}$.  It is
nontrivial, so Schur orthogonality gives zero for its character mean.  The
compact real form has an invariant positive-definite inner product, making
the adjoint representation self-dual and its character real.  Schur
orthogonality therefore also gives
$\int_G\chiad^2=\int_G|\chiad|^2=1$.  Finally, an orthogonal operator on the
$d_G$-dimensional real adjoint module has trace at most $d_G$, with equality
at the identity.
\end{proof}

\begin{lemma}[Even-moment lower bound for a positive cap]
\label{lem:moment-cap-lower-bound}
Let $X$ be real-valued with $-c\leq X\leq d$, let
$0<\varepsilon<d-c$, and put $t=d-\varepsilon$.  If $t\geq c$, then for
every integer $s\geq1$,
\[
 \mathbf P\{X>t\}
 \geq
 \frac{\mathbf E X^{2s}-t^{2s}}{d^{2s}-t^{2s}}.
\]
If moreover $\mathbf E X^2=1$ and $\varepsilon>1$, then
\[
 \mathbf P\{|X|<\varepsilon\}\geq1-\varepsilon^{-2}.
\]
\end{lemma}

\begin{proof}
On the complement of $\{X>t\}$ one has $|X|\leq t$, while everywhere
$|X|\leq d$.  Thus, with $p=\mathbf P\{X>t\}$,
\[
 \mathbf E X^{2s}\leq p d^{2s}+(1-p)t^{2s}.
\]
Rearrangement proves the first inequality.  The second is Markov's
inequality applied to $X^2$:
$\mathbf P\{|X|\geq\varepsilon\}\leq\varepsilon^{-2}\mathbf E X^2$.
\end{proof}

\begin{theorem}[Fixed-group total-degree tail]
\label{thm:fixed-group-total-degree-tail}
Let $G$ be a compact connected Lie group with simple Lie algebra, put
\[
 d_G=\dim\mathfrak g,
 \qquad
 c_G=-\min_{g\in G}\chiad(g),
\]
and choose
\[
 0<\varepsilon<\frac{d_G-2c_G}{2}.
\]
Define the positive Haar masses
\[
 p_{G,\varepsilon}
 =\mu_G\{g:\chiad(g)>d_G-\varepsilon\}\,
  \mu_G\{g:|\chiad(g)|<\varepsilon\}
\]
and put $q_{G,\varepsilon}=d_G-2\varepsilon$.  If $a\geq1$, $n$ is odd,
and
\[
 2a+n>
 \frac{\log(1/p_{G,\varepsilon})}
      {\log(q_{G,\varepsilon}/(2c_G))},
\]
then $I^G_{a,n}>0$.
\end{theorem}

\begin{proof}
First we verify that all constants in the statement exist.
Lemma~\ref{lem:adjoint-normalization} shows that the adjoint character has
Haar mean zero and second moment one, so it takes both positive and negative
values.  Since $G$ is connected, its continuous image
$\chiad(G)$ is an interval; hence that interval contains zero and $c_G>0$.
Garibaldi--Guralnick--Rains
\cite[Theorem~2.1]{GaribaldiGuralnickRains2025} prove
$\chiad(g)\geq-\operatorname{rank}(G)$.  If $r=\operatorname{rank}(G)$, then
\[
 d_G=r+|R|\geq3r>2r\geq2c_G,
\]
because a rank-$r$ root system contains the positive and negative simple
roots.  Thus the interval prescribed for $\varepsilon$ is nonempty and
$q_{G,\varepsilon}>2c_G$.  Both sets defining $p_{G,\varepsilon}$ are
nonempty open sets: the first contains the identity and the second contains
a point at which $\chiad$ vanishes.  Haar measure has full support on a
compact group: finitely many left translates of any nonempty open set cover
$G$, so such a set cannot have Haar measure zero.  Hence
$p_{G,\varepsilon}>0$.

Write $S=X+Y$ and $D=X-Y$.  On the negative half-plane $S<0$, the bounds
$X,Y\geq-c_G$ imply
\[
 |S|\leq2c_G,\qquad |D|\leq2c_G.
\]
Indeed, if $X\geq Y$, then $S<0$ gives $X<-Y\leq c_G$, whence
$0\leq X-Y<2c_G$; the other ordering is symmetric.  Therefore the absolute
value of the entire negative contribution is at most
\[
 (2c_G)^{2a+n}.
\]

On the positive rectangle
\[
 X>d_G-\varepsilon,\qquad |Y|<\varepsilon,
\]
one has $D>q_{G,\varepsilon}$ and $S>q_{G,\varepsilon}$.  Its contribution is
strictly greater than
\[
 p_{G,\varepsilon}q_{G,\varepsilon}^{\,2a+n}.
\]
The remainder of the half-plane $S\geq0$ is nonnegative because $n$ is odd
and $2a$ is even.  Consequently
\[
 I^G_{a,n}>
 (2c_G)^{2a+n}
 \left[
 p_{G,\varepsilon}
 \left(\frac{q_{G,\varepsilon}}{2c_G}\right)^{2a+n}-1
 \right],
\]
and the assumed logarithmic inequality makes the bracket positive.
\end{proof}

\begin{corollary}[Each fixed group has only a finite residual box]
\label{cor:fixed-group-finite-residual}
For a fixed compact connected $G$ with simple Lie algebra, only finitely many
pairs $(a,n)$ with $a\geq2$ and odd $n$ remain after combining the
two-minus Theorem~\ref{thm:two-minus-import} with
Theorem~\ref{thm:fixed-group-total-degree-tail}.
\end{corollary}

\begin{proof}
Theorem~\ref{thm:two-minus-import} supplies $a=1$.  The preceding theorem supplies every
pair with total degree $2a+n$ above one fixed finite threshold.  There are
only finitely many nonnegative integer pairs below that threshold.
\end{proof}

\section{Complete closure in type \texorpdfstring{$A_1$}{A1}}

\begin{lemma}[An effective $SU(2)$ Haar rectangle]
\label{lem:su2-effective-rectangle}
For the adjoint character of $\mathrm{SU}(2)$ one may take
\[
 c_G=1,\qquad d_G=3,\qquad \varepsilon=\frac18,
 \qquad q_{G,\varepsilon}=\frac{11}{4},
\]
and the Haar mass in Theorem~\ref{thm:fixed-group-total-degree-tail}
satisfies
\[
 p_{G,1/8}>
 \frac{3789107}{38106288000}>
 \left(\frac8{11}\right)^{29}.
\]
Consequently $I^{\mathrm{SU}(2)}_{a,n}>0$ whenever $a\geq1$, $n$ is odd,
and $2a+n\geq29$.
\end{lemma}

\begin{proof}
Parametrize conjugacy classes by $0\leq\alpha\leq\pi$.  The normalized
class density and adjoint character are
\[
 d\mu(\alpha)=\frac2\pi\sin^2\alpha\,d\alpha,
 \qquad
 \chiad(\alpha)=1+2\cos(2\alpha)=3-4\sin^2\alpha.
\]
Thus $c_G=1$ and $d_G=3$.  Put
\[
 A_0=(0,1/6)\cup(\pi-1/6,\pi).
\]
Since $\sin\alpha<\alpha$ on $(0,1/6)$, every point of $A_0$ satisfies
$\chiad>3-1/8$.  The elementary identity
\[
 \frac{22}{7}-\pi
 =\int_0^1\frac{x^4(1-x)^4}{1+x^2}\,dx>0
\]
gives $1/\pi>7/22$.  Taylor's theorem gives
$\sin x\geq x-x^3/6$ for $0\leq x\leq1/6$, and hence
\[
 \begin{aligned}
 \mu(A_0)
 &=\frac4\pi\int_0^{1/6}\sin^2x\,dx\\
 &>\frac{14}{11}\int_0^{1/6}(x-x^3/6)^2\,dx
 =\frac{541301}{277136640}.
 \end{aligned}
\]

Next put
\[
 B_0=\{\alpha:|\alpha-\pi/3|<1/32\}
 \cup\{\alpha:|\alpha-2\pi/3|<1/32\}.
\]
The character vanishes at the two interval centers and
$|\chiad'(\alpha)|\leq4$, so $|\chiad|<1/8$ on $B_0$.  For
$0\leq u\leq1/32$,
\[
 \sin(\pi/3-u)
 =\frac{\sqrt3}{2}\cos u-\frac12\sin u
 >\frac56\left(1-\frac1{2048}\right)-\frac1{64}
 =\frac{10043}{12288}>\frac45.
\]
Here we used $\sqrt3>5/3$, $\cos u\geq1-u^2/2$, and $\sin u\leq u$.
Symmetry gives the same bound on the second interval.  Their total length is
$1/8$, so
\[
 \mu(B_0)>
 \frac2\pi\frac18\frac{16}{25}>\frac{14}{275}.
\]
It follows that
\[
 p_{G,1/8}\geq\mu(A_0)\mu(B_0)
 >\frac{3789107}{38106288000}.
\]
The second displayed comparison in the statement follows from the exact
integer calculation
\[
 3789107\,11^{29}-38106288000\,8^{29}
 =114033204110077056951613191119358337>0.
\]
It is independently checked afresh by
\path{character_ring_iter/verify_full_q3_su2_gmp.cpp}.  Since
$q_{G,1/8}/(2c_G)=11/8$, the fixed-group tail theorem applies at every total
degree at least $29$.
\end{proof}

\begin{proposition}[Exact $SU(2)$ residual certificate]
\label{prop:su2-residual-certificate}
For every $a\geq2$ and every odd $n\geq1$ with
\[
 2a+n\leq27,
\]
one has $I^{\mathrm{SU}(2)}_{a,n}>0$.  There are exactly $78$ such pairs;
their smallest value is
\[
 I^{\mathrm{SU}(2)}_{2,1}=16.
\]
\end{proposition}

\begin{proof}
Let $c_k(j)$ be the multiplicity of the spin-$j$ irreducible in the $k$th
tensor power of the spin-one representation.  The $SU(2)$
Clebsch--Gordan rule gives the exact integer recurrence
\[
 c_0(j)=[j=0],\qquad
 c_{k+1}(0)=c_k(1),\qquad
 c_{k+1}(j)=c_k(j-1)+c_k(j)+c_k(j+1)\quad(j\geq1).
\]
The adjoint moment is $m_k=c_k(0)$.  Substitution in
Lemma~\ref{lem:full-moment-form} reduces every claimed inequality to a
finite integer sum.

The proof companion is
\begin{center}
\path{character_ring_iter/verify_full_q3_su2_gmp.cpp}.
\end{center}
It implements precisely this recurrence and the double sum in
Lemma~\ref{lem:full-moment-form}.  It uses GMP integers throughout, checks
every generated moment against the independent Catalan binomial transform
\[
 m_k=\sum_{j=0}^k(-1)^{k-j}\binom kj
 \frac1{j+1}\binom{2j}{j},
\]
which follows from $\chiad=|\operatorname{tr}U|^2-1$, and checks
the independently known prefix
\[
 (m_0,\ldots,m_9)=(1,0,1,1,3,6,15,36,91,232),
\]
enumerates the complete quantified domain, asserts that its cardinality is
$78$, rejects any nonpositive value, and asserts the stated exact minimum.
It also cross-multiplies the rational tail comparison in
Lemma~\ref{lem:su2-effective-rectangle}.  The command
\texttt{make full-q3-su2-audit} compiles the source with warnings promoted to
errors and reruns every check.  Its terminal success marker is
\texttt{SU2\_FULL\_Q3 VERIFICATION: ALL PASS}.
\end{proof}

\begin{theorem}[Full adjoint-generated $Q_3$ in type $A_1$]
\label{thm:su2-full-cone}
Let $G$ be any compact connected group with Lie algebra of type $A_1$.
Then Ginibre's full condition $Q_3$ holds on $\Qcone_G$.
\end{theorem}

\begin{proof}
By Lemma~\ref{lem:full-central-cover}, it suffices to use $G=\mathrm{SU}(2)$.
The automatic sectors are covered by
Theorem~\ref{thm:full-cone-reduction}.  The row $a=1$ is
Theorem~\ref{thm:two-minus-import}.  For $a\geq2$ and odd $n$, the total degree
$2a+n$ is odd.  Proposition~\ref{prop:su2-residual-certificate} covers
$2a+n\leq27$, while Lemma~\ref{lem:su2-effective-rectangle} covers
$2a+n\geq29$.  Hence the complete generator hierarchy is nonnegative, and
Theorem~\ref{thm:full-cone-reduction} promotes it to $\Qcone_G$.
\end{proof}

\section{Complete closure in type \texorpdfstring{$A_2$}{A2}}

\begin{proposition}[Exact $SU(3)$ cap and residual certificate]
\label{prop:su3-cap-residual-certificate}
For the adjoint character of $\mathrm{SU}(3)$,
\[
 m_{54}
 =1073719774632588649045995816896752254587680.
\]
This moment gives
\[
 \mathbf P\{X>6\}
 \geq\frac{m_{54}-6^{54}}{8^{54}-6^{54}},
 \qquad
 \mathbf P\{|X|<2\}\geq\frac34,
\]
and the exact comparison
\[
 \frac34\frac{m_{54}-6^{54}}{8^{54}-6^{54}}\cdot2^{29}>1
\]
holds.  Moreover,
\[
 I^{\mathrm{SU}(3)}_{a,n}>0
 \quad\text{for all }a\geq2,\ n\text{ odd},\ 2a+n\leq27.
\]
There are exactly $78$ inequalities in the last display, and their smallest
value is $I^{\mathrm{SU}(3)}_{2,1}=72$.
\end{proposition}

\begin{proof}
For Haar $U\in\mathrm{SU}(3)$, put $Z=|\operatorname{tr}U|^2$.  Scalar
invariance identifies its moments with the corresponding $\mathrm U(3)$
moments, and Schur--Weyl duality gives
\[
 E_j:=\mathbf E Z^j
 =\sum_{\substack{\lambda\vdash j\\\ell(\lambda)\leq3}}(f^\lambda)^2,
 \qquad
 m_k=\mathbf E(Z-1)^k
 =\sum_{j=0}^k(-1)^{k-j}\binom kjE_j,
\]
where $f^\lambda$ is computed by the hook-length formula.

The proof companion is
\begin{center}
\path{character_ring_iter/verify_full_q3_su3_gmp.cpp}.
\end{center}
It enumerates every partition in these formulas and performs all arithmetic with
GMP integers.  It reconstructs $m_0,\ldots,m_{54}$, checks the prefix
\[
 (m_0,\ldots,m_9)=(1,0,1,2,8,32,145,702,3598,19280),
\]
and independently verifies every generated moment against the proved
recurrence
\[
 (k+3)^2m_{k+1}=k(7k+11)m_k+8k(k+1)m_{k-1}\qquad(k\geq1).
\]
It asserts the displayed value of $m_{54}$ and cross-multiplies the tail
comparison exactly.

For the residual box it substitutes the reconstructed moments into
Lemma~\ref{lem:full-moment-form}, enumerates the complete $78$-pair domain,
rejects any nonpositive result, and asserts the exact minimum.  The command
\texttt{make full-q3-su3-audit} rebuilds and reruns the verifier.  Its
terminal marker is
\texttt{SU3\_FULL\_Q3 VERIFICATION: ALL PASS}.

It remains only to justify that the probability bounds consumed by the
integer comparison are valid.  The identity
$\chiad(U)=|\operatorname{tr}U|^2-1$ gives $-1\leq X\leq8$, and equality at
$-1$ occurs at $\operatorname{diag}(1,\omega,\omega^2)$ for a primitive cube
root $\omega$.  Apply Lemma~\ref{lem:moment-cap-lower-bound} with
$(c,d,\varepsilon,s)=(1,8,2,27)$.  Its two conclusions are precisely the
two displayed probability bounds.
\end{proof}

\begin{theorem}[Full adjoint-generated $Q_3$ in type $A_2$]
\label{thm:su3-full-cone}
Let $G$ be any compact connected group with Lie algebra of type $A_2$.
Then Ginibre's full condition $Q_3$ holds on $\Qcone_G$.
\end{theorem}

\begin{proof}
Central-cover invariance reduces the proof to $G=\mathrm{SU}(3)$.  The
automatic sectors and cone promotion are given by
Theorem~\ref{thm:full-cone-reduction}, and the $a=1$ row is
Theorem~\ref{thm:two-minus-import}.  Let $a\geq2$ and let $n$ be odd.  The character
bounds are $c_G=1,d_G=8$.  Taking $\varepsilon=2$ in
Theorem~\ref{thm:fixed-group-total-degree-tail} gives
$q_{G,\varepsilon}/(2c_G)=2$.  The exact probability comparison in
Proposition~\ref{prop:su3-cap-residual-certificate} therefore proves
$I_{a,n}>0$ for total degree $2a+n\geq29$.  The same proposition supplies
every odd total degree at most $27$.  This exhausts the hierarchy and proves
the theorem.
\end{proof}

\section{Complete closure in types \texorpdfstring{$A_3$ and $A_4$}{A3 and A4}}

\begin{proposition}[Exact $SU(4)$ and $SU(5)$ cap and residual certificates]
\label{prop:su45-cap-residual-certificate}
Put
\[
 M_4=m_{94}^{\mathrm{SU}(4)},\qquad
 M_5=m_{96}^{\mathrm{SU}(5)}.
\]
The following exact rational inequalities hold:
\[
 \frac{15}{16}
 \frac{M_4-11^{94}}{15^{94}-11^{94}}
 \left(\frac72\right)^{27}>1,
\]
and
\[
 \frac{285}{289}
 \frac{2^{96}M_5-31^{96}}{48^{96}-31^{96}}
 \left(\frac72\right)^{37}>1.
\]
In addition,
\[
 \begin{array}{c|c|c|c}
 G&\text{residual domain}&\text{number of pairs}&\text{minimum}\\
 \hline
 \mathrm{SU}(4)&a\geq2,\ n\text{ odd},\ 2a+n\leq25
   &66&I_{2,1}=94,\\
 \mathrm{SU}(5)&a\geq2,\ n\text{ odd},\ 2a+n\leq35
   &136&I_{2,1}=96.
 \end{array}
\]
Every residual value in the table is strictly positive.
\end{proposition}

\begin{proof}
For $N=4,5$, Schur--Weyl duality and the hook-length formula give the exact
source
\[
 E_j^{(N)}=
 \sum_{\substack{\lambda\vdash j\\\ell(\lambda)\leq N}}(f^\lambda)^2,
 \qquad
 m_k^{(N)}=\sum_{j=0}^k(-1)^{k-j}\binom kjE_j^{(N)}.
\]
The proof companion is
\begin{center}
\path{character_ring_iter/verify_full_q3_su45_gmp.cpp}.
\end{center}
It enumerates these partitions, evaluates every hook product and binomial
transform with GMP integers, and reconstructs the moments through degrees
$94$ and $96$.  As independent checks, it verifies the prefixes
\[
 \begin{aligned}
 (m_0^{(4)},\ldots,m_9^{(4)})
   &=(1,0,1,2,9,43,245,1557,10829,80958),\\
 (m_0^{(5)},\ldots,m_9^{(5)})
   &=(1,0,1,2,9,44,264,1824,14210,122224),
 \end{aligned}
\]
and every instance in range of the two proved low-rank recurrences from
Parts~I--II.  It then substitutes the moments in
Lemma~\ref{lem:full-moment-form}, checks the complete residual domains and
their exact minima, and cross-multiplies both displayed tail inequalities.
No floating-point comparison is used.  The command
\texttt{make full-q3-su45-audit} rebuilds the verifier; its terminal marker is
\texttt{SU45\_FULL\_Q3 VERIFICATION: ALL PASS}.

To connect the exact comparisons to Haar probabilities, use
$\chiad=|\operatorname{tr}U|^2-1$, so $c_G=1$ and
$d_G=N^2-1$.  Trace-zero elements exist in both groups.  For $SU(4)$ apply
Lemma~\ref{lem:moment-cap-lower-bound} with
\[
 (d,\varepsilon,s)=(15,4,47).
\]
It gives cap mass at least $(M_4-11^{94})/(15^{94}-11^{94})$ and central
mass at least $15/16$; the tail ratio is
$(d-2\varepsilon)/(2c_G)=7/2$.  For $SU(5)$ use
\[
 (d,\varepsilon,s)=(24,17/2,48).
\]
After multiplying numerator and denominator by $2^{96}$, its cap bound is
$(2^{96}M_5-31^{96})/(48^{96}-31^{96})$; the central mass is at least
$285/289$, and the tail ratio is again $7/2$.  These are exactly the
quantities checked by the verifier.
\end{proof}

\begin{theorem}[Full adjoint-generated $Q_3$ in types $A_3$ and $A_4$]
\label{thm:su45-full-cone}
Let $G$ be a compact connected group whose simple Lie algebra has type
$A_3$ or $A_4$.  Then Ginibre's full condition $Q_3$ holds on $\Qcone_G$.
\end{theorem}

\begin{proof}
Central-cover invariance reduces the two cases to $SU(4)$ and $SU(5)$.
The automatic sectors and cone promotion follow from
Theorem~\ref{thm:full-cone-reduction}; the $a=1$ row is supplied by
Theorem~\ref{thm:two-minus-import}.  For $SU(4)$, the first exact comparison of
Proposition~\ref{prop:su45-cap-residual-certificate} and
Theorem~\ref{thm:fixed-group-total-degree-tail} supply every odd total
degree at least $27$, while the residual certificate supplies every odd
total degree at most $25$.  The corresponding split for $SU(5)$ is
$2a+n\geq37$ and $2a+n\leq35$.  Thus no generator case remains in either
group, and Theorem~\ref{thm:full-cone-reduction} completes the proof.
\end{proof}

\section{Stable type \texorpdfstring{$A$}{A}}

\begin{theorem}[Full stable-rank type-$A$ hierarchy]
\label{thm:stable-a-full}
Let $G=\mathrm{SU}(N)$ and let $a\geq1$, $n\geq0$.  If
\[
 N\geq 2a+n,
\]
then $I^{G}_{a,n}>0$.  In fact,
\[
 F_a^G(t)=(2a)!\,\frac{e^{-2t}}{(1-t)^{2a+2}}
 \quad\text{through degree }n.
\]
\end{theorem}

\begin{proof}
For every $r\leq N$, Schur--Weyl duality gives
\[
 \mathbf E_{\mathrm{SU}(N)}|\operatorname{tr}U|^{2r}=r!.
\]
Since $\chiad(U)=|\operatorname{tr}U|^2-1$, all adjoint moments through
degree $N$ agree with those of $A=E-1$, where $E$ is exponential with mean
one.  The integral $I_{a,n}$ uses no moment above degree $2a+n$, so the
rank hypothesis permits replacement of $X,Y$ by independent copies of
$A$.

Let $E_1,E_2$ be independent mean-one exponentials, put
$R=E_1+E_2$ and $T=E_1/R$.  Then $R$ and $T$ are independent,
$R$ has the gamma density $re^{-r}\,dr$, and $T$ is uniform on $[0,1]$.
Moreover,
\[
 X-Y=R(2T-1),\qquad X+Y=R-2.
\]
Consequently
\[
 \begin{aligned}
 F_a^G(t)
 &=\mathbf E(2T-1)^{2a}\,
   \mathbf E\bigl[R^{2a}e^{t(R-2)}\bigr]\\
 &=\frac1{2a+1}e^{-2t}
   \frac{(2a+1)!}{(1-t)^{2a+2}}
  =(2a)!\frac{e^{-2t}}{(1-t)^{2a+2}}.
 \end{aligned}
\]
The logarithm of the last factor after removal of the positive constant is
\[
 -2t-(2a+2)\log(1-t).
\]
Its first cumulant is $2a$, and its $r$th cumulant for $r\geq2$ is
$(2a+2)(r-1)!$.  They are all positive, so
Lemma~\ref{lem:nonnegative-cumulants} proves strict positivity of every
coefficient.
\end{proof}

\begin{lemma}[Paley--Zygmund cap bound]\label{lem:paley-cap}
If $W\geq0$ has finite second moment and $0\leq t<\mathbf EW$, then
\[
 \mathbf P\{W>t\}\geq
 \frac{(\mathbf EW-t)^2}{\mathbf E W^2}.
\]
\end{lemma}

\begin{proof}
Write $A=\{W>t\}$.  Then
\[
 \mathbf EW
 \leq t+\mathbf E(W\mathbf1_A)
 \leq t+(\mathbf EW^2)^{1/2}\mathbf P(A)^{1/2}
\]
by Cauchy--Schwarz.  Rearrangement and squaring prove the claim.
\end{proof}

\begin{lemma}[Polynomial one-sided tail bound]
\label{lem:polynomial-one-sided-tail}
Let $X$ be real-valued, let $T\in\mathbb R$, and let $P$ be a real
polynomial for which
\[
 A=\mathbf E\bigl[(X-T)P(X)^2\bigr]>0,
 \qquad
 B=\mathbf E\bigl[(X-T)^2P(X)^4\bigr]>0.
\]
Then
\[
 \mathbf P\{X>T\}\geq\frac{A^2}{B}.
\]
\end{lemma}

\begin{proof}
Put
\[
 R(x)=1-\frac AB(x-T)P(x)^2.
\]
If $x\leq T$, then $R(x)\geq1$, so
$\mathbf1_{\{x\leq T\}}\leq R(x)^2$.  Therefore
\[
 \begin{aligned}
 \mathbf P\{X\leq T\}
 &\leq \mathbf E R(X)^2\\
 &=1-2\frac AB A+\frac{A^2}{B^2}B
 =1-\frac{A^2}{B}.
 \end{aligned}
\]
Taking complements proves the assertion.
\end{proof}

\begin{lemma}[One-sided cap propagation for every even-minus sector]
\label{lem:one-sided-cap-full-q3}
Let $X,Y$ be independent copies of a real random variable satisfying
$X\geq-C$ and $\mathbf EX^2=1$, and let $T,\delta,p_0$ satisfy
\[
 \delta>1,\qquad
 \mathbf P\{X>T\}\geq p_0>0,\qquad
 q:=T-\delta>2C.
\]
If $K$ is odd and
\begin{equation}
\label{eq:one-sided-cap-full-q3-test}
 p_0(1-\delta^{-2})q^K>(2C)^K,
\end{equation}
then
\[
 \mathbf E[(X-Y)^{2a}(X+Y)^n]>0
\]
for every $a\geq1$, every odd $n$, and every $2a+n\geq K$.
\end{lemma}

\begin{proof}
Chebyshev's inequality and $\mathbf EY^2=1$ give
$\mathbf P\{|Y|<\delta\}\geq1-\delta^{-2}$.  On the product of this
event with $\{X>T\}$, both $X-Y$ and $X+Y$ exceed $q$.  On the negative
half-plane $X+Y<0$, the lower support bound gives
$|X-Y|,|X+Y|\leq2C$, as in the proof of
Theorem~\ref{thm:fixed-group-total-degree-tail}.  The rest of the positive
half-plane contributes nonnegatively.  Hence, for $L=2a+n$,
\[
 \mathbf E[(X-Y)^{2a}(X+Y)^n]
 >p_0(1-\delta^{-2})q^L-(2C)^L.
\]
The right side is positive at $L=K$ by
\eqref{eq:one-sided-cap-full-q3-test}; after division by $(2C)^L$ it is
strictly increasing in $L$ because $q/(2C)>1$.
\end{proof}

\begin{proposition}[Uniform higher-minus tail for $SU(N)$, $N\geq6$]
\label{prop:sun-ge6-tail}
Let $a\geq1$ and let $n$ be odd.  Put $L=2a+n$ and define
\[
 K_N=
 \begin{cases}
 27,&N=6,\\
 29,&N=7,8,\\
 31,&N\geq9.
 \end{cases}
\]
If $N\geq6$ and $L\geq K_N$, then
\[
 I^{\mathrm{SU}(N)}_{a,n}>0.
\]
\end{proposition}

\begin{proof}
Put $Z=|\operatorname{tr}U|^2$ and $X=Z-1$.  For
\[
 E_r^{(N)}=\mathbf E Z^r
 =\sum_{\substack{\lambda\vdash r\\\ell(\lambda)\leq N}}(f^\lambda)^2,
\]
apply Lemma~\ref{lem:paley-cap} to $W=Z^{16}$ and
$t=(33/5)^{16}$.  The exact inequalities verified below include
$E_{16}^{(N)}>t$ and give
\[
 \mathbf P\{X>28/5\}
 =\mathbf P\{Z>33/5\}
 \geq
 \frac{(E_{16}^{(N)}-(33/5)^{16})^2}{E_{32}^{(N)}}.
\]
Since $\mathbf EX^2=1$, Lemma~\ref{lem:moment-cap-lower-bound} also gives
\[
 \mathbf P\{|X|<11/10\}\geq\frac{21}{121}.
\]
On the product of these two events, both $X-Y$ and $X+Y$ exceed
\[
 \frac{28}{5}-\frac{11}{10}=\frac92.
\]
On the negative half-plane $X+Y<0$, the identity $X=Z-1\geq-1$ gives
$|X-Y|,|X+Y|\leq2$.  Exactly as in
Theorem~\ref{thm:fixed-group-total-degree-tail}, it therefore suffices that
\[
 \frac{21}{121}
 \frac{(E_{16}^{(N)}-(33/5)^{16})^2}{E_{32}^{(N)}}
 \left(\frac94\right)^L>1.
\]

The exact GMP verifier
\begin{center}
\path{character_ring_iter/verify_full_q3_sun_ge6_gmp.cpp}
\end{center}
clears denominators and checks this inequality at $L=K_N$ for every
$6\leq N\leq31$.  For $N\geq32$, Schur--Weyl stability gives
$E_{16}^{(N)}=16!$ and $E_{32}^{(N)}=32!$, so the single checked stable-row
comparison at $N=32$ applies to every larger $N$.  The left side increases
with $L$, because $9/4>1$, completing the tail proof.
\end{proof}

\begin{proposition}[Exact residual type-$A$ certificate for $N\geq6$]
\label{prop:sun-ge6-residual}
For every $6\leq N\leq28$ and every pair $a\geq2$, $n$ odd, satisfying
\[
 N<2a+n<K_N,
\]
one has $I^{\mathrm{SU}(N)}_{a,n}>0$.  The quantified domain contains
$1315$ pairs.  Its smallest value is
\[
 I^{\mathrm{SU}(6)}_{2,3}=3550.
\]
\end{proposition}

\begin{proof}
The proof companion is
\begin{center}
\path{character_ring_iter/verify_full_q3_sun_ge6_gmp.cpp}.
\end{center}
For each $6\leq N\leq28$, it reconstructs
$E_0^{(N)},\ldots,E_{29}^{(N)}$ by enumerating all partitions of bounded
length and evaluating the hook-length formula, then forms
\[
 m_k^{(N)}=\sum_{j=0}^k(-1)^{k-j}\binom kjE_j^{(N)}.
\]
It checks $E_j^{(N)}=j!$ independently whenever $j\leq N$, substitutes the
moments in Lemma~\ref{lem:full-moment-form}, enforces the exact stable/tail
complement defining the residual domain, checks the per-rank and total scope
counts, rejects every nonpositive value, and asserts the displayed minimum.

The same source reconstructs $E_{16}^{(N)}$ and $E_{32}^{(N)}$ for
$6\leq N\leq31$, uses $16!,32!$ at the stable endpoint, and checks the
tail comparison of Proposition~\ref{prop:sun-ge6-tail} after the exact
integer clearing
\[
 21\bigl(5^{16}E_{16}^{(N)}-33^{16}\bigr)^2 9^{K_N}
 >121\,5^{32}E_{32}^{(N)}4^{K_N}.
\]
The command \texttt{make full-q3-sun-ge6-audit} rebuilds and reruns the
complete calculation.  Its terminal marker is
\texttt{SUN\_GE6\_FULL\_Q3 VERIFICATION: ALL PASS}.
\end{proof}

\begin{theorem}[Full adjoint-generated $Q_3$ in every type $A$]
\label{thm:type-a-full-cone}
Let $G$ be a compact connected group whose simple Lie algebra has type
$A_r$ for some $r\geq1$.  Then Ginibre's full condition $Q_3$ holds on
$\Qcone_G$.
\end{theorem}

\begin{proof}
Types $A_1$--$A_4$ are Theorems~\ref{thm:su2-full-cone},
\ref{thm:su3-full-cone}, and~\ref{thm:su45-full-cone}.  Let
$G=SU(N)$ with $N\geq6$; central-cover invariance handles every other global
form.  The automatic sectors follow from
Theorem~\ref{thm:full-cone-reduction}, and
Theorem~\ref{thm:two-minus-import} supplies $a=1$.

It remains to take $a\geq2$, $n$ odd, and $L=2a+n$.  If $L\leq N$,
Theorem~\ref{thm:stable-a-full} applies.  If $L\geq K_N$,
Proposition~\ref{prop:sun-ge6-tail} applies.  A pair in neither range can
occur only for $6\leq N\leq28$, and every such pair is supplied by
Proposition~\ref{prop:sun-ge6-residual}.  Indeed, for $N=29$ there is no odd
integer strictly between $29$ and $31$, and for $N\geq30$ the stable and
tail ranges overlap.  Thus the generator hierarchy is complete for every
$N$.  Theorem~\ref{thm:full-cone-reduction} promotes it to $\Qcone_G$.
\end{proof}

\section{Complete closure for the exceptional groups}

For an exceptional group $G$, write $m_j=m_j^G$ and set
\[
 A_G=m_{2k+1}-T m_{2k},\qquad
 B_G=m_{4k+2}-2T m_{4k+1}+T^2m_{4k}.
\]
The following table fixes all parameters used below.  The column $M$ is the
largest moment consumed by the tail calculation, and $K$ is an odd total
degree.
\begin{table}[ht]
\centering
\begin{tabular}{c|rrrrrr}
$G$&$c$&$T$&$\delta$&$k$&$M$&$K$\\ \hline
$G_2$&2&91/10&6/5&9&38&25\\
$F_4$&4&211/10&13/10&12&50&37\\
$E_6$&3&98/5&7/5&10&42&29\\
$E_7$&7&333/10&13/10&17&70&63\\
$E_8$&8&461/10&13/10&23&94&71
\end{tabular}
\caption{Parameters for the exceptional higher-minus tail certificates.}
\label{tab:exceptional-full-tail}
\end{table}

\begin{proposition}[Exact exceptional higher-minus tails]
\label{prop:exceptional-full-tail}
Take the row of Table~\ref{tab:exceptional-full-tail} corresponding to $G$.
If $a\geq1$, $n$ is odd, and $2a+n\geq K$, then
\[
 I^G_{a,n}>0.
\]
\end{proposition}

\begin{proof}
The bounds $\chiad\geq-c$ used in the five rows are rigorous: for
$G_2,F_4,E_7,E_8$ they follow from the rank bound of
Garibaldi--Guralnick--Rains~\cite[Theorem~2.1]{GaribaldiGuralnickRains2025},
and for $E_6$ the sharper value $\min\chiad=-3$ is
\cite[Theorem~2.3 and Table~1]{GaribaldiGuralnickRains2025}.

Apply Lemma~\ref{lem:polynomial-one-sided-tail} with $P(x)=x^k$.  It gives
\[
 \mathbf P\{X>T\}\geq \frac{A_G^2}{B_G}.
\]
By Lemma~\ref{lem:adjoint-normalization} and Markov's inequality,
\[
 \mathbf P\{|Y|<\delta\}\geq1-\delta^{-2}.
\]
On the product of these two events,
\[
 X-Y>T-\delta,
 \qquad
 X+Y>T-\delta.
\]
On $X+Y<0$, the common lower bound $X,Y\geq-c$ instead gives
$|X-Y|,|X+Y|\leq2c$, exactly as in
Theorem~\ref{thm:fixed-group-total-degree-tail}.  Thus, with $L=2a+n$, it
suffices that
\begin{equation}
\label{eq:exceptional-full-tail-test}
 \rho_G R_G^L>1,
 \qquad
 \rho_G=\frac{A_G^2}{B_G}(1-\delta^{-2}),
 \qquad
 R_G=\frac{T-\delta}{2c}.
\end{equation}
The exact verifier
\begin{center}
\path{character_ring_iter/verify_full_q3_exceptional_gmp.cpp}
\end{center}
checks $A_G,B_G>0$, $R_G>1$, and
$\rho_GR_G^K>1$ by GMP rational arithmetic in every row.  Hence
\eqref{eq:exceptional-full-tail-test} holds for every $L\geq K$.
\end{proof}

\begin{proposition}[Exact exceptional residual certificate]
\label{prop:exceptional-full-residual}
For each row of Table~\ref{tab:exceptional-full-tail}, one has
\[
 I^G_{a,n}>0
 \quad\text{whenever}\quad
 a\geq2,\quad n\text{ is odd},\quad 2a+n<K.
\]
The exact scopes and minima are
\[
\begin{array}{c|rrr}
G&\text{largest residual degree}&\text{number of pairs}&
 \min I^G_{a,n}\\ \hline
G_2&23&55&36\\
F_4&35&136&36\\
E_6&27&78&38\\
E_7&61&435&36\\
E_8&69&561&36
\end{array}
\]
and every displayed minimum occurs at $(a,n)=(2,1)$.
\end{proposition}

\begin{proof}
The proof companion is
\begin{center}
\path{character_ring_iter/verify_full_q3_exceptional_gmp.cpp}.
\end{center}
It reads the exact invariant-tensor moments already regenerated from the
Cartan data by the source audit of Parts~I--II.  For $E_8$, the independently
regenerated values through degree $70$ are joined to the transcribed degrees
$71$--$100$ from the ancillary tensor-rank table of
Bourjaily--Plesser--Vergu
\cite[ancillary file \texttt{adjoint\_tensor\_ranks.tex}]{BPV2024}; the
Parts~I--II source replay compares that transcription against its own
Racah--Speiser reconstruction through $\chiad^{\otimes50}$.

For every pair in the stated boxes, the verifier evaluates the integer
formula of Lemma~\ref{lem:full-moment-form}, rejects a nonpositive value,
checks the five scope counts, and asserts the five minima and their
coordinates.  Independently, it forms $A_G,B_G,\rho_G,R_G$ as reduced GMP
rationals, checks the tail inequality at $K$, and checks that the same chosen
certificate does not yet pass at $K-2$.  The command
\texttt{make full-q3-exceptional-audit} rebuilds and reruns the arithmetic
audit; its terminal marker is
\texttt{EXCEPTIONAL\_FULL\_Q3 VERIFICATION: ALL PASS}.
\end{proof}

\begin{theorem}[Full adjoint-generated $Q_3$ for every exceptional type]
\label{thm:exceptional-full-cone}
Let $G$ be a compact connected group whose simple Lie algebra has type
$G_2,F_4,E_6,E_7$, or $E_8$.  Then Ginibre's full condition $Q_3$ holds on
$\Qcone_G$.
\end{theorem}

\begin{proof}
By central-cover invariance, it suffices to use the compact form employed in
the moment computation.  The automatic sectors follow from
Theorem~\ref{thm:full-cone-reduction}, and
Theorem~\ref{thm:two-minus-import} supplies $a=1$.  For
$a\geq2$ and odd $n$, Proposition~\ref{prop:exceptional-full-tail} applies
when $2a+n\geq K$, while
Proposition~\ref{prop:exceptional-full-residual} applies below $K$.  Thus the
entire generator hierarchy holds, and
Theorem~\ref{thm:full-cone-reduction} promotes it to $\Qcone_G$.
\end{proof}

\section{Stable classical types}

\begin{theorem}[Full stable classical hierarchy]
\label{thm:stable-bcd-full}
Let $a\geq1$ and $n\geq0$, and put $L=2a+n$.  The inequality
$I^G_{a,n}\geq0$ holds in each of the following ranges:
\[
 \begin{array}{c|c}
 G&\text{stable-rank condition}\\
 \hline
 B_\ell=\mathrm{SO}(2\ell+1)&L\leq2\ell+1,\\
 C_\ell=\mathrm{Sp}(2\ell)&L\leq2\ell+1,\\
 D_\ell=\mathrm{SO}(2\ell)&L\leq\ell-1.
 \end{array}
\]
More precisely, throughout these ranges
\[
 F_a^G(t)=e^{-t+t^2/2}
 \sum_{j=0}^{a}\frac{\gamma_{a,j}}{(1-t)^{2j+1}}
 \quad\text{through degree }n,
\]
where
\[
 \gamma_{a,j}
 =\binom{2a}{2j}(2a-2j-1)!!
   \frac{\binom{2j}{j}}{4^j}(2j)!>0.
\]
Here $(-1)!!=1$.
\end{theorem}

\begin{proof}
The exact stable joint trace laws of Diaconis--Shahshahani
\cite[Theorems~4 and~6]{DiaconisShahshahani1994} identify the stable
orthogonal and symplectic adjoint variables with the common law
\[
 X=\frac{Z_1^2-1+\sqrt2Z_2}{2},
\]
where $Z_1,Z_2$ are independent standard normal variables.  The sign of
$Z_2$ is immaterial.

We first check that the stated ranks supply every moment used by
$I_{a,n}$.  For type $B$, the orthogonal Schur integral criterion
\cite[Section~3, paragraph preceding Lemma~3.1]{Rains1997} applied to
$s_{(1,1)}^s$ is stable when $s\leq2\ell+1$.  For type $C$, the
symplectic criterion in the same paragraph is stable when
$\ell\geq\lfloor s/2\rfloor$, equivalently $s\leq2\ell+1$.  In type $D$,
the trace degree of $\chiad^s$ is $2s$; the exact orthogonal trace law and
the equality of the $\mathrm{SO}(2\ell)$ and $\mathrm O(2\ell)$ integrals
apply when $2s<2\ell$, namely $s\leq\ell-1$.  Since
Lemma~\ref{lem:full-moment-form} uses only moments of total degree at most
$L$, the displayed hypotheses give the stable law throughout the required
calculation.

Take independent copies
\[
 X=\frac{Z_1^2-1+\sqrt2Z_2}{2},\qquad
 Y=\frac{W_1^2-1+\sqrt2W_2}{2}.
\]
In polar coordinates for $(Z_1,W_1)$, put
\[
 R=\frac{Z_1^2+W_1^2}{2},\qquad
 C=\frac{Z_1^2-W_1^2}{Z_1^2+W_1^2}.
\]
Then $R$ is mean-one exponential, $C=\cos(2\Theta)$ with $\Theta$ uniform,
and $R,C$ are independent.  Also
\[
 A=\frac{Z_2+W_2}{\sqrt2},\qquad
 B=\frac{Z_2-W_2}{\sqrt2}
\]
are independent standard normals, independent of $(R,C)$, and
\[
 X+Y=R-1+A,\qquad X-Y=RC+B.
\]
Only even powers survive in the conditional expansion of $(RC+B)^{2a}$.
Using
\[
 \mathbf E B^{2q}=(2q-1)!!,
 \qquad
 \mathbf E C^{2j}=\frac{\binom{2j}{j}}{4^j},
 \qquad
 \mathbf E[R^{2j}e^{tR}]=\frac{(2j)!}{(1-t)^{2j+1}},
\]
and $\mathbf E e^{tA}=e^{t^2/2}$ gives the asserted formula for
$F_a^G$.

It remains to prove coefficientwise positivity.  For each $j\geq0$, put
\[
 H_j(t)=\frac{e^{-t+t^2/2}}{(1-t)^{2j+1}}.
\]
Its logarithm has cumulants
\[
 \kappa_1=2j,\qquad
 \kappa_2=2j+2,\qquad
 \kappa_r=(2j+1)(r-1)!\quad(r\geq3).
\]
They are nonnegative, so every coefficient of $H_j$ is nonnegative by
Lemma~\ref{lem:nonnegative-cumulants}.  The coefficients $\gamma_{a,j}$
are strictly positive; hence their finite linear combination has
nonnegative coefficients as claimed.
\end{proof}

\section{A simultaneous nonstable classical tail}

For the rest of the classical argument let $X=\chiad(U)$ and
$Y=\chiad(V)$ for independent Haar elements, and put
\[
 S=X+Y,\qquad D=X-Y.
\]
For $G=B_b,C_b,D_b$, use the defining trace $T_1$ and the signed square
trace
\[
 \tau_G=
 \begin{cases}
  \operatorname{Tr}(U^2),&G=B_b,D_b,\\
  -\operatorname{Tr}(U^2),&G=C_b.
 \end{cases}
\]
Thus in all three families
\begin{equation}
\label{eq:classical-adjoint-trace-form}
 X=\frac{T_1^2-\tau_G}{2}.
\end{equation}

\begin{lemma}[Exact classical negative endpoints]
\label{lem:full-classical-negative-endpoints}
For the standard representatives,
\[
 C_{B_b}=b,\qquad C_{C_b}=b,\qquad
 C_{D_b}=\begin{cases}b,&b\text{ even},\\ b-2,&b\text{ odd}.
 \end{cases}
\]
\end{lemma}

\begin{proof}
Write the defining eigenvalue pairs as $z_j,z_j^{-1}$, put
$x_j=(z_j+z_j^{-1})/2\in[-1,1]$, and set
$R=\sum_jx_j$ and $Q=\sum_jx_j^2$.  Direct substitution in
\eqref{eq:classical-adjoint-trace-form} gives
\[
 \chiad^{B_b}=b+2R+2(R^2-Q),\qquad
 \chiad^{D_b}=b+2(R^2-Q),\qquad
 \chiad^{C_b}=2R^2+2Q-b.
\]
The first two expressions are affine in each variable separately, so their
minima occur at vertices of the cube.  At a vertex $Q=b$ and
$R=b-2k$.  The type-$B$ value is $2R^2+2R-b$, whose minimum is $-b$;
the type-$D$ value is $2R^2-b$, whose minimum is $-b$ for even $b$ and
$2-b$ for odd $b$.  The type-$C$ expression is at least $-b$, with
equality at $x_1=\cdots=x_b=0$.
\end{proof}

\begin{lemma}[Root-normalized classical Weyl cap]
\label{lem:full-q3-classical-weyl-cap}
Let $G=B_b$ or $D_b$, put $d=\dim G$, and use the positive roots
\[
 \begin{array}{c|c|c|c}
 G&\kappa\text{ in }Q(\theta)=\sum_{\alpha>0}\alpha(\theta)^2
   =\kappa|\theta|^2&|W|&(d_1,\ldots,d_b)\\ \hline
 B_b&2b-1&2^bb!&(2,4,\ldots,2b)\\
 D_b&2b-2&2^{b-1}b!&(2,4,\ldots,2b-2,b).
 \end{array}
\]
Set $z_{B_b}=2^{-b}$ and $z_{D_b}=1$.  Suppose that $Q(\theta)\leq R$
lies in the standard eigenangle chart and that throughout this ball
\[
 \prod_{\alpha>0}
 \left(\frac{\sin(\alpha(\theta)/2)}{\alpha(\theta)/2}\right)^2
 \geq\eta>0.
\]
Then, for the adjoint character $X$,
\begin{equation}
\label{eq:full-q3-classical-cap-mass}
 \mathbf P\{X>d-R\}\geq
 \frac{z_G\eta}{|W|}
 \frac{\prod_{i=1}^b d_i!}
 {(2\pi)^{b/2}\Gamma(d/2+1)}
 \left(\frac{R}{2\kappa}\right)^{d/2}.
\end{equation}
If $0<R<24$, one may take
\begin{equation}
\label{eq:full-q3-product-sine-small}
 \eta=\left(1-\frac R{24}\right)^2.
\end{equation}
For $D_b$, one may instead take, whenever $R<12\kappa$,
\begin{equation}
\label{eq:full-q3-product-sine-long}
 \eta=\left(1-\frac{R}{12\kappa}\right)^\kappa.
\end{equation}
\end{lemma}

\begin{proof}
The character estimate $2(1-\cos u)\leq u^2$ gives
$d-X\leq Q(\theta)$.  If $Q(\theta)=R>0$, at least one root coordinate is
nonzero and $2(1-\cos u)<u^2$ for that coordinate; if $Q(\theta)<R$, the
conclusion is immediate.  Hence $X>d-R$ throughout the cap.  Weyl
integration and the stated sine-product bound
therefore reduce the cap mass to
\[
 \frac{z_G\eta}{|W|(2\pi)^b}
 \int_{Q(\theta)\leq R}\prod_{\alpha>0}\alpha(\theta)^2\,d\theta.
\]
The factor $z_{B_b}=2^{-b}$ occurs because the squared-length-two
Macdonald--Mehta normalization replaces each short root $e_i$ by
$\sqrt2e_i$; all type-$D$ roots already have squared length two.  The
Macdonald--Mehta identity
\cite[Theorem~3.1(i)]{Etingof2009}, followed by radial integration of the
homogeneous discriminant, gives
\[
 \int_{Q(\theta)\leq R}\prod_{\alpha>0}\alpha(\theta)^2\,d\theta
 =(2\pi)^{b/2}\frac{\prod_i d_i!}{\Gamma(d/2+1)}
 \left(\frac{R}{2\kappa}\right)^{d/2},
\]
which proves~\eqref{eq:full-q3-classical-cap-mass}.

For $u=\alpha(\theta)$, the elementary bound
$\sin(u/2)/(u/2)\geq1-u^2/24$ applies.  If $R<24$, all factors on the
right are positive, and
\[
 \prod_{\alpha>0}(1-u_\alpha)^2
 \geq\left(1-\sum_{\alpha>0}u_\alpha\right)^2,
 \qquad u_\alpha=\frac{\alpha(\theta)^2}{24},
\]
proves~\eqref{eq:full-q3-product-sine-small}.  In type $D$, every root has
squared length two, so
$u_\alpha\leq R/(12\kappa)$.  Concavity of $\log(1-u)$ on the interval
from zero to this upper bound, together with
$\sum_\alpha u_\alpha=Q(\theta)/24\leq R/24$, gives
\eqref{eq:full-q3-product-sine-long}.
\end{proof}

\begin{lemma}[All-even-minus trace-tail reduction]
\label{lem:all-minus-trace-tail}
Let $C_G=-\min_G\chiad$, and choose $c>2C_G$, $d>1$, and
$q\in\mathbb R$.  Define
\[
 A=\{|T_1|\geq\sqrt{2c+2d+q},\ \tau_G\leq q\},
 \qquad B=\{|X|<d\},
 \qquad p=\mathbf P(A)\mathbf P(B).
\]
If $a\geq1$, $n$ is odd, and $L=2a+n$, then
\begin{equation}
\label{eq:all-minus-base-tail}
 I^G_{a,n}\geq 2pc^L-2\mathbf E(\tau_G)_+^L.
\end{equation}
More generally, if $L_0\leq L$, then
\begin{equation}
\label{eq:all-minus-propagated-tail}
 I^G_{a,n}\geq
 2pc^L-2(2C_G)^{L-L_0}\mathbf E(\tau_G)_+^{L_0}.
\end{equation}
Consequently, if $L_0$ is odd and
\begin{equation}
\label{eq:all-minus-tail-test}
 pc^{L_0}>\mathbf E(\tau_G)_+^{L_0},
\end{equation}
then $I^G_{a,n}>0$ for every $a\geq1$, odd $n$, and
$2a+n\geq L_0$.
\end{lemma}

\begin{proof}
On $A$, equation~\eqref{eq:classical-adjoint-trace-form} gives
$X\geq c+d$.  Hence on the rectangle $A\times B$ both $S$ and $D$ are at
least $c$.  On the transposed rectangle $B\times A$, one has
$S\geq c$ and $D\leq-c$.  These rectangles are disjoint because
$X\geq c+d>d$ on $A$, and each has probability $p$.  Their combined
contribution is therefore at least $2pc^L$.

The rest of the half-plane $S\geq0$ contributes nonnegatively.  On
$S<0$, split into $X\leq Y$ and $Y<X$.  The contributions have equal
absolute value.  On the first piece $X<0$ and
\[
 |D|=Y-X\leq-2X,
 \qquad
 |S|=-X-Y\leq-2X.
\]
Equation~\eqref{eq:classical-adjoint-trace-form} also gives
$-2X\leq\tau_G$, so $|D|,|S|\leq(\tau_G)_+$.  This proves
\eqref{eq:all-minus-base-tail} after doubling.  Independently,
$X,Y\geq-C_G$ implies $|D|,|S|\leq2C_G$ on $S<0$.  Retaining $L_0$ of
the $L$ factors under the first bound and the other $L-L_0$ factors under
the second proves~\eqref{eq:all-minus-propagated-tail}.

Finally, if~\eqref{eq:all-minus-tail-test} holds, then $c>2C_G$ makes the
right side of~\eqref{eq:all-minus-propagated-tail} strictly positive for
every $L\geq L_0$.  Every relevant total degree is odd, which proves the
last assertion.
\end{proof}

\begin{lemma}[Exponential bound for the square-trace moment]
\label{lem:square-trace-moment-mgf}
For every integer $L\geq1$ and every $v>0$,
\[
 \mathbf E(\tau_G)_+^L
 \leq\left(\frac{L}{ev}\right)^L\mathbf E e^{v\tau_G}.
\]
For the classical square traces and $C_G\geq296$, one has in particular
\begin{equation}
\label{eq:square-trace-v4}
 \mathbf E(\tau_G)_+^L
 \leq4e^{20}\left(\frac{L}{4e}\right)^L.
\end{equation}
\end{lemma}

\begin{proof}
For $z\geq0$, the maximum of $z^Le^{-vz}$ occurs at $z=L/v$ and equals
$(L/(ev))^L$.  Multiplication by $e^{vz}$ and expectation prove the first
display.

For $B_b$, the power-two identity of Rains
\cite[Theorem~1.5, equation~(23)]{RainsPowerMaps} gives
\[
 \tau_G\mathrel{\mathop=\limits^{\mathrm d}}
 1+2\operatorname{Re}\operatorname{Tr}W,
 \qquad W\sim\Haar(\mathrm U(b)).
\]
The strong Szeg\H{o} upper bound of
Courteaut--Johansson--Lambert
\cite[Lemma~2.12]{CJL2022} therefore gives
$\mathbf E e^{v\tau_G}\leq e^{v+v^2}$.

For $C_b,D_b$, the corresponding power-two decomposition
\cite[Theorem~1.5, equation~(22), and equation~(36)]{RainsPowerMaps}
writes $\tau_G-1$ as a sum, up to harmless sign changes, of two independent
centered full orthogonal traces.  The Fredholm estimate obtained from
\cite[Proposition~3.3 and equations~(3.4)--(3.6)]{CJL2022} is
\[
 \mathbf E e^{v\tau_G}
 \leq
 \frac{e^{v+v^2}}
 {(1-\mathfrak b_{d_0}(v))(1-\mathfrak b_{d_1}(v))},
 \qquad
 \mathfrak b_d(v)=
 \frac{e^{v^2/(d+1)}}{1-v^2/(d+1)^2}\frac{v^d}{d!},
\]
provided the companion trace-norm envelope is below one.  When
$C_G\geq296$, both component dimensions are at least $11$.  At $v=4$,
the directed verifier cited below checks
\[
 \mathfrak a_d(v)=
 \frac{e^{v^2/(d+1)}v^d}
 {d!(1-v/(d+1))\sqrt{1-v^2/(d+1)^2}},
 \qquad
 \mathfrak a_{11}(4)<1,
 \qquad \mathfrak b_{11}(4)<\frac12.
\]
Both quantities decrease with $d\geq11$: after replacing $d$ by $d+1$,
the factorial contributes $4/(d+1)$, while the exponential and each
remaining rational factor also decrease.  Thus each denominator is at
least $1/2$.  In every family
$\mathbf E e^{4\tau_G}\leq4e^{20}$, and substitution in the first display
proves~\eqref{eq:square-trace-v4}.
\end{proof}

\begin{lemma}[Source-audited high-rank trace rectangle]
\label{lem:full-q3-high-rank-trace-rectangle}
Put
\[
 r=\frac{2000001}{1000000},\qquad
 \eta=\frac{359}{2000},\qquad
 A_0=\frac{4571}{2000}.
\]
For $G=B_b,C_b,D_b$ with $C=C_G\geq296$, set
\[
 c=rC,\qquad d=q=\eta C.
\]
Then
\[
 \mathbf P\{|T_1|\geq\sqrt{2c+3d},\ \tau_G\leq d\}
 \geq\frac12e^{-A_0C},
 \qquad
 \mathbf P\{|X|<d\}\geq\frac34.
\]
\end{lemma}

\begin{proof}
The second inequality is Chebyshev's inequality, since
$\mathbf EX^2=1$ and $d>2$.

The first is the Rains-square trace certificate of Parts~I--II, whose
inputs are recalled here.  The one-trace density estimate
\cite[Theorem~1.4]{CJL2022}, the Mills lower bound, and the checked
orthogonal parity shifts give
\[
 \mathbf P\{|T_1|\geq\sqrt{(2r+3\eta)C}\}
 \geq e^{-A_0C}.
\]
Rains' power-two decomposition and the same density estimate give
\[
 \mathbf P\{\tau_G\geq\eta C\}
 \leq
 \exp\!\left(-\frac{(\eta C-1)^2}{4}\right)
 +\frac{10(\log(2\lfloor C/2\rfloor))^{1/4}}
 {\sqrt{\Gamma(2\lfloor C/2\rfloor+1)}}
 \leq\frac12e^{-A_0C}.
\]
The last two displayed comparisons, including the unitary projection in
type $B$, are checked with directed rounding for every
$296\leq C\leq10000$ and with explicit derivative signs thereafter by
\path{character_ring_iter/trace_square_cutoff_mpfr.cpp}.  Subtracting the
square-trace exceptional event proves the result.
\end{proof}

\begin{proposition}[Full higher-minus tail for large classical rank]
\label{prop:bcd-high-rank-full-tail}
Let $G=B_b,C_b,D_b$ and suppose $C_G\geq296$.  Then
\[
 I^G_{a,n}>0
 \qquad(a\geq1,\ n\text{ odd})
\]
outside the stable range of Theorem~\ref{thm:stable-bcd-full}.  Hence the
complete generator hierarchy holds in every such row.
\end{proposition}

\begin{proof}
Let $L_0$ be the first odd total degree after the stable range.  The exact
negative endpoints give
\[
 L_0=
 \begin{cases}
  2b+3=2C_G+3,&G=B_b,C_b,\\
  b+1=C_G+1,&G=D_b,\ b\text{ even},\\
  b=C_G+2,&G=D_b,\ b\text{ odd}.
 \end{cases}
\]
In particular $C_G+1\leq L_0\leq2C_G+3$.

Use the parameters of
Lemma~\ref{lem:full-q3-high-rank-trace-rectangle}.  In the notation of
Lemma~\ref{lem:all-minus-trace-tail}, its two probability estimates give
\[
 p\geq\frac38e^{-A_0C},\qquad C=C_G.
\]
Lemma~\ref{lem:square-trace-moment-mgf} shows that the ratio of the two
sides in~\eqref{eq:all-minus-tail-test} is at least
\[
 \frac3{32}e^{-A_0C-20}
 \left(\frac{4erC}{L_0}\right)^{L_0}.
\]
Since the base is greater than one and
$C+1\leq L_0\leq2C+3$, the logarithm of this ratio is bounded below by
\begin{equation}
\label{eq:bcd-high-rank-margin}
 H(C)=\log\frac3{32}-A_0C-20
 +(C+1)\log\left(\frac{4erC}{2C+3}\right).
\end{equation}
The directed-rounding companion
\begin{center}
\path{character_ring_iter/verify_full_q3_bcd_high_mpfr.cpp}
\end{center}
checks $H(296)>8$, $H'(296)>1/10$, and, more importantly for the
monotonicity argument,
\[
 \log\left(\frac{4er\cdot296}{2\cdot296+3}\right)-A_0
 >\frac9{100}.
\]
The factor $4erC/(2C+3)$ is increasing, and the remaining term in
\[
 H'(C)=-A_0+\log\left(\frac{4erC}{2C+3}\right)
 +(C+1)\left(\frac1C-\frac2{2C+3}\right)
\]
is positive.  Consequently the first two terms in $H'(C)$ have sum
greater than $9/100$ for every $C\geq296$, while the last term is positive.
Hence $H'(C)>0$ throughout that range.  Thus
\eqref{eq:all-minus-tail-test} holds at $L_0$, and
Lemma~\ref{lem:all-minus-trace-tail} proves every later odd total degree.
Theorem~\ref{thm:stable-bcd-full} supplies the prefix.

The same verifier checks the two Fredholm values consumed in
Lemma~\ref{lem:square-trace-moment-mgf}.  The command
\texttt{make full-q3-bcd-high-audit} rebuilds both directed-rounding
programs and requires their terminal zero-failure markers.  The
source-matched machine-C transcript
\path{certificates/full_q3/fullq3bcdanalytic0003_machine_c.log}
authenticates both sources and records zero compile and run statuses; its
adjacent SHA-256 file authenticates the transcript.
\end{proof}

\begin{lemma}[Finite-rank Fredholm MGF bounds]
\label{lem:finite-rank-fredholm-mgf}
Let
\[
 d_T=\begin{cases}2b+1,&G=B_b,\\2b,&G=C_b,D_b,
 \end{cases}
 \qquad
 \mathfrak b_d(v)=
 \frac{e^{v^2/(d+1)}}{1-v^2/(d+1)^2}\frac{v^d}{d!},
\]
and put
\[
 \mathfrak a_d(v)=
 \frac{e^{v^2/(d+1)}v^d}
 {d!(1-v/(d+1))\sqrt{1-v^2/(d+1)^2}},
 \qquad D_d(v)=1-\mathfrak b_d(v).
\]
If
$0<v<d_T+1$ and $\mathfrak a_{d_T}(v)<1$, then
\begin{equation}
\label{eq:finite-trace-mgf-sandwich}
 e^{v^2/2}D_{d_T}(v)
 \leq\mathbf E e^{\pm vT_1}
 \leq\frac{e^{v^2/2}}{D_{d_T}(v)}.
\end{equation}
For $C_b$ put
\[
 d_0=2\left\lceil\frac{b+1}{2}\right\rceil,
 \qquad d_1=2\left\lceil\frac b2\right\rceil+1,
\]
and for $D_b$ put
\[
 d_0=2\left\lceil\frac b2\right\rceil,
 \qquad d_1=2\left\lceil\frac{b-1}{2}\right\rceil+1.
\]
Whenever the corresponding two trace-norm conditions hold,
\begin{equation}
\label{eq:finite-square-mgf-upper}
 \mathbf E e^{v\tau_G}\leq
 \begin{cases}
 e^{v+v^2},&G=B_b,\\[2pt]
 \displaystyle\frac{e^{v+v^2}}{D_{d_0}(v)D_{d_1}(v)},
   &G=C_b,D_b.
 \end{cases}
\end{equation}
\end{lemma}

\begin{proof}
We spell out the Fredholm step because every finite-middle trace estimate
uses it.  Courteaut--Johansson--Lambert
\cite[Proposition~3.3]{CJL2022} state their Fredholm identity for every
complex \(\xi\), not only for real characteristic-function parameters.
Substitution of \(\xi=\mp iv\) therefore writes the MGF of the
random-eigenvalue trace as a Gaussian exponential times
\[
 \det(I+QK_{\pm v}Q).
\]
Here \(Q\) deletes the finite prefix determined by the Jacobi ensemble and
\(K_{\pm v}\) is the corresponding Bessel Hankel operator.  In the proof of
\cite[Lemma~3.6, equation~(3.4)]{CJL2022}, the row-norm inequality followed
by the Bessel bound and the two geometric sums gives, before the source
imposes its simpler sufficient restriction \(\lvert z\rvert<2e^{-5/4}n\),
the explicit estimate
\[
 \lVert QK_{\pm v}Q\rVert_1
 \leq \mathfrak a_{d_T}(v).
\]
Thus the hypothesis \(\mathfrak a_{d_T}(v)<1\) justifies the absolutely
convergent Plemelj series
\[
 \log\det(I+QK_{\pm v}Q)
 =\sum_{j\geq1}\frac{(-1)^{j+1}}j
   \operatorname{tr}(QK_{\pm v}Q)^j.
\]
The source's total matrix size \(d\) is the component dimension denoted by
\(d_T\) here, and its projection \(Q_n\) starts at the corresponding Jacobi
index.  The argument of
\cite[Lemma~3.7, equation~(3.6)]{CJL2022}, again before replacing the
explicit expression by a simpler sufficient range, gives
\[
 \sum_{j\geq1}\frac1j
  \left|\operatorname{tr}(QK_{\pm v}Q)^j\right|
 \leq-\log\bigl(1-\mathfrak b_{d_T}(v)\bigr).
\]
Notice that \(0\leq\mathfrak b_d(v)\leq\mathfrak a_d(v)<1\): after writing
\(x=v/(d+1)\), the ratio \(\mathfrak a_d/\mathfrak b_d\) is
\(\sqrt{(1+x)/(1-x)}\).  Exponentiating the last absolute bound therefore
gives
\[
 D_{d_T}(v)
 \leq\det(I+QK_{\pm v}Q)
 \leq D_{d_T}(v)^{-1}.
\]
The determinant is positive because the Gaussian prefactor is positive and
the left side of the Fredholm identity is an MGF.  This proves the determinant
sandwich directly from the cited estimates.

It remains to account for the deterministic eigenvalues.  The source formula
is written for the random eigenvalues only.  In type $B$, their random trace
has mean $-1$; the fixed $+1$ eigenvalue cancels that shift exactly.  Equivalently,
the source factor \(e^{\mp v}\) cancels the restoration factor \(e^{\pm v}\)
in the two MGFs.  In types~$C,D$ there is no deterministic shift.  The full
defining character is therefore centered in every family, in agreement with
character orthogonality, and the Gaussian prefactor is \(e^{v^2/2}\).  This
proves~\eqref{eq:finite-trace-mgf-sandwich} for both signs.

For the square trace, Rains' power-two decompositions cited in
Lemma~\ref{lem:square-trace-moment-mgf} give either the type-$B$ unitary trace
or two independent orthogonal components of dimensions \(d_0,d_1\) displayed
in the statement.  Apply the preceding upper determinant bound to each
orthogonal component and multiply, using independence.  This gives the
\(C,D\) clause of~\eqref{eq:finite-square-mgf-upper}; the type-$B$ clause is
the unitary strong Szeg\H{o} bound already proved in that lemma.
\end{proof}

\begin{lemma}[Exact defining-trace determinant formulas]
\label{lem:exact-defining-trace-determinants}
For $s\geq0$, write $I_k=I_k(2s)$ for the modified Bessel function.  If
$T_1$ is the full defining trace, then
\begin{align}
 \mathbf E_{B_b}e^{sT_1}
 &=e^s\det[I_{|j-k|}-I_{j+k+1}]_{j,k=0}^{b-1},
 \label{eq:exact-b-positive-mgf}\\
 \mathbf E_{B_b}e^{-sT_1}
 &=e^{-s}\det[I_{|j-k|}+I_{j+k+1}]_{j,k=0}^{b-1},
 \label{eq:exact-b-negative-mgf}\\
 \mathbf E_{C_b}e^{sT_1}
 &=\det[I_{|j-k|}-I_{j+k+2}]_{j,k=0}^{b-1},
 \label{eq:exact-c-mgf}\\
 \mathbf E_{D_b}e^{sT_1}
 &=\frac12\det[I_{|j-k|}+I_{j+k}]_{j,k=0}^{b-1}.
 \label{eq:exact-d-mgf}
\end{align}
The type-$C$ and type-$D$ trace laws are symmetric.
\end{lemma}

\begin{proof}
Courteaut--Johansson--Lambert~\cite[Lemma~3.1]{CJL2022} give the four
Toeplitz--Hankel determinants for the Jacobi parameters in their
equation~(3.1).  Substitution of $\psi(x)=e^{2sx}$ gives Fourier
coefficients $I_k(2s)$.  Their $E_n^{+-}$ case is
$\mathrm{SO}(2b+1)$ after restoring the fixed $+1$ eigenvalue, giving
\eqref{eq:exact-b-positive-mgf}; replacing $s$ by $-s$ gives
\eqref{eq:exact-b-negative-mgf}.  Their $E_n^{++}$ case simultaneously
describes $\mathrm{Sp}(2b)$ and the random part of
$\mathrm O^-(2b+2)$, giving~\eqref{eq:exact-c-mgf}.  Their $E_n^{--}$
case is $\mathrm{SO}(2b)$; its determinant equals $2$ at $s=0$, which
accounts for the factor $1/2$ in~\eqref{eq:exact-d-mgf}.  Finally,
$-I$ belongs to the type-$C$ and even-orthogonal groups, proving symmetry.
\end{proof}

\begin{lemma}[Finite layered recovery from an exact MGF]
\label{lem:finite-layered-mgf-recovery}
Let $Z$ be bounded, let $M(s)=\mathbf Ee^{sZ}$, and fix $s>0$ and a lower
endpoint $t$.  Choose $\lambda_->0$, increasing upper endpoints
$u_0<\cdots<u_r$, and positive $\lambda_0,\ldots,\lambda_r$.  Put
\[
 L_-=e^{\lambda_-t}\frac{M(s-\lambda_-)}{M(s)},
 \qquad
 U_i=e^{-\lambda_i u_i}\frac{M(s+\lambda_i)}{M(s)},
 \qquad F_i=1-L_--U_i,
 \qquad w_i=e^{-su_i}.
\]
If all $F_i$ are positive, then
\begin{equation}
\label{eq:finite-layered-mgf-recovery}
 \mathbf P\{Z>t\}\geq M(s)
 \left[w_rF_r+\sum_{i=0}^{r-1}(w_i-w_{i+1})F_i\right].
\end{equation}
An endpoint at or above the upper support of $Z$ may be used with $U_i=0$.
\end{lemma}

\begin{proof}
Under the tilted probability
$d\mathbf P_s=e^{sZ}M(s)^{-1}d\mathbf P$, exponential Markov gives
\[
 \mathbf P_s\{Z\leq t\}\leq L_-,
 \qquad
 \mathbf P_s\{Z\geq u_i\}\leq U_i.
\]
Thus $\mathbf P_s\{t<Z<u_i\}\geq F_i$.  On the other hand, summation by
parts over the layers
\[
 (t,u_0),\ [u_0,u_1),\ \ldots,\ [u_{r-1},u_r)
\]
gives
\[
 \mathbf E_s[e^{-sZ}\mathbf1_{\{t<Z<u_r\}}]
 \geq w_r\mathbf P_s\{t<Z<u_r\}
 +\sum_{i=0}^{r-1}(w_i-w_{i+1})
  \mathbf P_s\{t<Z<u_i\}.
\]
Insert the preceding lower bounds and multiply by $M(s)$.
\end{proof}

\begin{lemma}[One-step tilted defining-trace lower bound]
\label{lem:one-step-trace-tilt}
Let $t,h>0$ and put
\[
 s=t+h,
 \qquad
 \rho(t,h)=1-
 \frac{e^{-h^2/2}}{D(s)}
 \left(\frac1{D(s-h)}
       +\frac1{D(s+h)}\right),
\]
where $D(v)=D_{d_T}(v)$.  Assume $s-h>0$, the
sandwich~\eqref{eq:finite-trace-mgf-sandwich} at
$s-h,s,s+h$, and $\rho(t,h)>0$.  Then
\begin{equation}
\label{eq:one-step-trace-tilt}
 \mathbf P\{|T_1|\geq t\}
 \geq2D(s)\rho(t,h)
 \exp\left\{\frac{s^2}{2}-s(t+2h)\right\}.
\end{equation}
\end{lemma}

\begin{proof}
For $\sigma\in\{+1,-1\}$, tilt probability by
$M_\sigma(s)^{-1}e^{\sigma sT_1}$.  Exponential Markov at parameter $h$
and the two sides of~\eqref{eq:finite-trace-mgf-sandwich} show that the
tilted probability of leaving the interval
$t<\sigma T_1<t+2h$ is at most $1-\rho(t,h)$.  The choice
$s=t+h$ centers the Gaussian comparison at the midpoint of this interval;
both exit losses are exactly $e^{-h^2/2}$.  On the interval the inverse tilt
is at least $e^{-s(t+2h)}$.  The lower MGF bound at $s$ therefore
gives
\[
 \mathbf P\{\sigma T_1\geq t\}
 \geq D(s)\rho(t,h)e^{s^2/2-s(t+2h)}.
\]
The positive and negative tails are disjoint.  Summing the two bounds proves
the display; no symmetry of the trace law is being assumed.
\end{proof}

\begin{proposition}[Full higher-minus tail for the finite middle ranks]
\label{prop:bcd-middle-rank-full-tail}
The complete generator hierarchy holds in the following classical rows:
\[
 B_b\ (22\leq b\leq295),\qquad
 C_b\ (29\leq b\leq295),\qquad
 D_b\ (71\leq b\leq297).
\]
\end{proposition}

\begin{proof}
Let $L_0$ be the first odd degree after the stable range, as in the proof of
Proposition~\ref{prop:bcd-high-rank-full-tail}.  In every row put
\[
 r=\frac{2000001}{1000000},\qquad c=rC_G,
 \qquad d=\frac32,
 \qquad t=\sqrt{2c+2d+q},
\]
and use the family-dependent parameters
\[
\begin{array}{c|c|c|c}
G&\text{rank range}&q&h\\ \hline
B_b&22\leq b\leq295&\frac{21}{5}\sqrt b&\frac{13}{10}\\
C_b&29\leq b\leq295&4\sqrt b&\frac{13}{10}\\
D_b&71\leq b\leq297&\frac{18}{5}\sqrt b&\frac65.
\end{array}
\]
Let $V_G$ denote the right side of
\eqref{eq:one-step-trace-tilt}.  At
\[
 v_q=\frac{q-1}{2},
 \]
exponential Markov and~\eqref{eq:finite-square-mgf-upper} give an upper
bound $U_G$ for $\mathbf P\{\tau_G\geq q\}$.  Hence the event $A$ in
Lemma~\ref{lem:all-minus-trace-tail} has probability at least $V_G-U_G$.
Chebyshev gives $\mathbf P(B)\geq5/9$.  Finally, at
\[
 v_m=\frac{\sqrt{1+8L_0}-1}{4},
\]
Lemmas~\ref{lem:square-trace-moment-mgf} and
\ref{lem:finite-rank-fredholm-mgf} bound
$\mathbf E(\tau_G)_+^{L_0}$.

The directed-rounding companion
\begin{center}
\path{character_ring_iter/verify_full_q3_bcd_mid_mpfr.cpp}
\end{center}
 evaluates precisely these formulas.  It checks every trace-norm domain,
the positive $\rho$ value for both one-sided tilts, the strict difference
$V_G-U_G$, and
\[
 \frac59(V_G-U_G)c^{L_0}>
 \mathbf E(\tau_G)_+^{L_0}
\]
in all $768$ rows.  Its minimum directed logarithmic margin is
\[
 1.0047268114794421860\ldots>0
 \quad\text{at }D_{71};
\]
the minimum log separation between $V_G$ and $U_G$ is
$0.16429898256055109\ldots$ at $C_{29}$, and the largest trace-norm
envelope is $0.39958077942628398\ldots<1$, also at $C_{29}$.
Thus~\eqref{eq:all-minus-tail-test} holds at $L_0$ in every row.
Lemma~\ref{lem:all-minus-trace-tail} supplies the entire tail and
Theorem~\ref{thm:stable-bcd-full} supplies the prefix.
\end{proof}

\begin{proposition}[Additional finite classical tails]
\label{prop:bcd-low-tail}
For the rows $B_{18},\ldots,B_{21}$, let the odd tail onsets be
\[
 (K_{18},K_{19},K_{20},K_{21})=(47,45,45,45).
\]
For $D_{31},\ldots,D_{70}$, in increasing rank order, let the onsets be
\[
\begin{split}
(&49,51,51,51,51,53,53,55,53,55,
55,57,55,57,57,59,59,61,61,61,\\
&61,63,61,63,63,65,63,65,65,67,
67,69,67,69,69,71,71,71,71,73).
\end{split}
\]
Then $I^G_{a,n}>0$ whenever $a\geq1$, $n$ is odd, and
$2a+n\geq K_b$ in the corresponding row.
\end{proposition}

\begin{proof}
Use $c=(2000001/1000000)C_G$ in every row.  For $B_{18}$ take
\[
 d=\frac{13}{10},\qquad q=\frac65b,\qquad h=\frac54;
\]
for $B_{19},B_{20},B_{21}$ take
\[
 d=\frac32,\qquad q=b,\qquad h=\frac65.
\]
For the type-$D$ rows take $d=3/2$, $h=6/5$, and
\[
 q=\begin{cases}
  7b/10,&31\leq b\leq38,\\
  3b/5,&39\leq b\leq52,\\
  b/2,&53\leq b\leq70.
 \end{cases}
\]
The probability and moment bounds are exactly those in the proof of
Proposition~\ref{prop:bcd-middle-rank-full-tail}, now evaluated at the
displayed onset rather than at the first post-stable degree.  The same
directed-rounding verifier checks all $44$ rows.  Its least logarithmic
tail margin is
\[
 0.07529217823349327949\ldots>0
 \quad\text{at }D_{43},\ K_{43}=55;
\]
its least tilted mass parameter is
$0.02269453416254314484\ldots>0$ at $B_{19}$.  Its least logarithmic
separation between the defining-trace lower bound and square-trace upper
bound, attained at $D_{40}$, is
\[
 0.83270005200917997982\ldots>0.
\]
Consequently~\eqref{eq:all-minus-tail-test} holds at every displayed
onset.  Lemma~\ref{lem:all-minus-trace-tail} propagates each inequality to
all later odd total degrees.
\end{proof}

\begin{lemma}[Exact finite classical adjoint moments]
\label{lem:finite-orthogonal-adjoint-moments}
Write
\[
 e_2^j=\sum_{\lambda\vdash2j}c_j(\lambda)s_\lambda,
 \qquad c_j(\lambda)\in\mathbb Z_{\geq0},
\]
and let $s_j$ be the stable classical adjoint moment.  Then
\[
 m_j^{B_b}=s_j-
 \sum_{\substack{\nu\vdash j\\\ell(\nu)>2b+1}}
 c_j(2\nu).
\]
For type $C$ one has
\[
 m_j^{C_b}=s_j-
 \sum_{\substack{\nu\vdash j\\\ell(\nu)>b}}
 c_j\!\left((\nu_1,\nu_1,\nu_2,\nu_2,\ldots)'\right).
\]
For type $D$ one has, without a restriction on $j$,
\[
 m_j^{D_b}=s_j-
 \sum_{\substack{\nu\vdash j\\\ell(\nu)>2b}}c_j(2\nu)
 +\sum_{\substack{\nu\vdash j-b\\\ell(\nu)\leq2b}}
 c_j\bigl(\lambda(b,\nu)\bigr),
\]
where, after padding $\nu$ by zeros to length $2b$,
\[
 \lambda(b,\nu)=(1+2\nu_1,\ldots,1+2\nu_{2b}).
\]
An empty partition sum is zero.
\end{lemma}

\begin{proof}
For both orthogonal families the adjoint representation is
$\bigwedge^2V$, so the displayed Schur expansion is the character
decomposition of $(\bigwedge^2V)^{\otimes j}$; iterated vertical-two-strip
Pieri proves $c_j(\lambda)\in\mathbb Z_{\geq0}$.  The orthogonal Schur
integral criterion in Rains~\cite[Section~3, paragraph preceding
Lemma~3.1]{Rains1997} says that the $\mathrm O(N)$ integral of $s_\lambda$
is one precisely for even $\lambda$ of length at most $N$, and is zero
otherwise.  For odd $N=2b+1$, the central element $-I$ identifies the
$\mathrm O(N)$ and $\mathrm{SO}(N)$ integrals of every power of
$\bigwedge^2V$.  Subtracting the even partitions excluded by the length
wall gives the type-$B$ formula.

For type $C$, the adjoint representation is $\operatorname{Sym}^2V$.
The standard involution $\omega$ of the symmetric-function ring sends
$h_2$ to $e_2$ and $s_\lambda$ to $s_{\lambda'}$.  Thus the coefficient of
$s_\lambda$ in $h_2^j$ is $c_j(\lambda')$.  Rains' symplectic criterion in
the same paragraph says that the integral of $s_\lambda$ over
$\mathrm{Sp}(2b)$ is one precisely when $\ell(\lambda)\leq2b$ and every
row length of $\lambda$ occurs an even number of times, and is zero
otherwise.  Such a partition is uniquely
$\lambda=(\nu_1,\nu_1,\nu_2,\nu_2,\ldots)$; its length constraint is
$\ell(\nu)\leq b$.  Subtraction from the unrestricted stable sum proves the
type-$C$ formula.

For $\mathrm{SO}(2b)$, averaging a polynomial character over the identity
component equals its $\mathrm O(2b)$ average plus its determinant-twisted
$\mathrm O(2b)$ average.  The first term admits the even partitions of
length at most $2b$.  The second admits precisely the partitions of length
$2b$ whose rows are all odd: subtracting one full column leaves an even
partition, to which Rains' orthogonal criterion applies.  Every such
determinant-class partition is uniquely $\lambda(b,\nu)$, and its size
equation is
$2b+2|\nu|=2j$.  Hence $|\nu|=j-b$, proving the type-$D$ formula.
\end{proof}

\begin{lemma}[Bounded Littlewood determinants for the finite classical moments]
\label{lem:bounded-littlewood-classical-moments}
Work in the completed ring of symmetric functions over $\mathbb Q$, put
$e_r=h_r=0$ for $r<0$, and define
\[
 f_r^e=\sum_{k\in\mathbb Z}e_ke_{k+r},\qquad
 f_r^h=\sum_{k\in\mathbb Z}h_kh_{k+r},\qquad
 e=\sum_{k\geq0}e_k,\qquad
 \bar e=\sum_{k\geq0}(-1)^ke_k.
\]
For $b\geq1$, let
\[
\begin{aligned}
 \mathcal B_b
 &=\frac12\left\{
 e\det_{1\leq i,j\leq b}(f^e_{i-j}-f^e_{i+j-1})
 +\bar e\det_{1\leq i,j\leq b}(f^e_{i-j}+f^e_{i+j-1})
 \right\},\\
 \mathcal C_b
 &=\det_{1\leq i,j\leq b}(f^h_{i-j}-f^h_{i+j}),\\
 \mathcal D_b
 &=\frac12\det_{1\leq i,j\leq b}(f^e_{i-j}+f^e_{i+j-2}).
\end{aligned}
\]
Then, for every $j\geq0$,
\[
 m_j^{B_b}=[m_{(2^j)}]\mathcal B_b,\qquad
 m_j^{C_b}=[m_{(2^j)}]\mathcal C_b,\qquad
 m_j^{D_b}=[m_{(2^j)}]\mathcal D_b.
\]
Here $[m_{(2^j)}]$ denotes the coefficient of the monomial symmetric
function $m_{(2^j)}$, and every displayed completed series is interpreted
degree by degree.
\end{lemma}

\begin{proof}
Let $c(\lambda)$ and $r(\lambda)$ denote the numbers of odd columns and odd
rows of $\lambda$.  Hall duality gives
\[
 \langle h_2^j,F\rangle=[m_{(2^j)}]F.
\]
Transpose the admitted partitions in the type-$B$ formula of
Lemma~\ref{lem:finite-orthogonal-adjoint-moments}.  Since the involution
$\omega$ interchanges $e_2$ and $h_2$, this gives
\[
 m_j^{B_b}
 =\left\langle h_2^j,
   \sum_{\substack{\lambda_1\leq2b+1\\c(\lambda)=0}}s_\lambda
  \right\rangle.
\]
Huh--Kim--Krattenthaler--Okada
\cite[Theorem~1.3, equations~(1.9)--(1.10)]{HKKO2025}, specialized at
$u=0$, give respectively the sum of the $c(\lambda)=0$ and
$c(\lambda)=2b+1$ sectors and their difference.  Taking their half-sum gives
exactly $\mathcal B_b$.

For type $C$, transposition gives
\[
 m_j^{C_b}
 =\left\langle e_2^j,
   \sum_{\substack{\lambda_1\leq2b\\r(\lambda)=0}}s_\lambda
  \right\rangle.
\]
Equation~(1.6) of the same source identifies the Schur sum with
$\det(f^e_{i-j}-f^e_{i+j})$.  Apply the Hall-isometric involution $\omega$;
it sends $e_2$ to $h_2$ and $f_r^e$ to $f_r^h$.  Hall duality then gives the
formula for $\mathcal C_b$.

Finally, the two components in the type-$D$ formula transpose to the sectors
\[
 \lambda_1\leq2b,\qquad c(\lambda)=0\ \text{or}\ c(\lambda)=2b.
\]
The second sector is precisely the determinant-twisted component: its $2b$
columns all have odd length.  Equation~(1.7) of the cited theorem at $u=0$
is the sum of these two sectors and equals $\mathcal D_b$.  A final use of
Hall duality proves the type-$D$ identity.
\end{proof}

\begin{lemma}[Two-power coefficient functional]
\label{lem:two-power-coefficient-functional}
Let $F$ be homogeneous of symmetric-function degree $2j$.  Define the
specialization $\pi_2$ on power sums by
\[
 \pi_2(p_1)=p,\qquad \pi_2(p_2)=q,\qquad
 \pi_2(p_r)=0\quad(r\geq3).
\]
Then
\begin{equation}
\label{eq:two-power-coefficient-functional}
 [m_{(2^j)}]F
 =j!\sum_{a=0}^{j}(2(j-a)-1)!!
 [p^{2(j-a)}q^a]\pi_2(F).
\end{equation}
Under this specialization,
\[
 \sum_{r\geq0}\pi_2(e_r)u^r
 =\exp(pu-qu^2/2),\qquad
 \sum_{r\geq0}\pi_2(h_r)u^r
 =\exp(pu+qu^2/2),
\]
and hence, with $e_{-1}=h_{-1}=0$,
\[
 r e_r=pe_{r-1}-qe_{r-2},\qquad
 r h_r=ph_{r-1}+qh_{r-2}.
\]
In the last display $e_r$ and $h_r$ abbreviate their specialized images
$\pi_2(e_r)$ and $\pi_2(h_r)$.
\end{lemma}

\begin{proof}
Expand $F$ in the power-sum basis.  A power-sum monomial containing $p_r$
with $r\geq3$ cannot contain the ordinary monomial
$x_1^2\cdots x_j^2$, so it contributes nothing to
$[m_{(2^j)}]F$.  The remaining power-sum monomials are
$p_1^{2(j-a)}p_2^a$.  In their expansion, the $a$ ordered $p_2$ factors must
choose distinct variables, while the remaining $2(j-a)$ ordered $p_1$
factors must hit each of the other variables twice.  Thus the coefficient of
$x_1^2\cdots x_j^2$ is
\[
 \frac{j!}{(j-a)!}\frac{(2(j-a))!}{2^{j-a}}
 =j!(2(j-a)-1)!!,
\]
which proves~\eqref{eq:two-power-coefficient-functional}.  The two generating
series are the standard power-sum formulas for elementary and complete
symmetric functions after applying $\pi_2$.  Differentiating them with
respect to $u$ and comparing coefficients gives the recurrences.
\end{proof}

\begin{proposition}[Exact polynomial and Weyl-cap tails in the residual rows]
\label{prop:bcd-remaining-exact-tails}
For each row in the following table, $I^G_{a,n}>0$ whenever $a\geq1$,
$n$ is odd, and $2a+n\geq K$:
\[
\begin{array}{c|rrrr|c}
G&T&\delta&k&K&\text{method}\\ \hline
B_3&227/20&23/20&8&37&\text{polynomial}\\
B_4&331/20&5/4&10&41&\text{polynomial}\\
B_5&75/4&6/5&10&55&\text{polynomial}\\
C_3&57/5&6/5&8&37&\text{polynomial}\\
D_4&57/4&23/20&9&47&\text{polynomial}\\
D_5&131/10&6/5&7&33&\text{polynomial}
\end{array}
\qquad
\begin{array}{c|rrrr}
G&R&\delta&q&K\\ \hline
B_2,C_2&-&6/5&39/5&47\\
B_6&1053/50&3/2&1386/25&51\\
D_6&371/20&7/5&921/20&49\\
D_7&18&8/5&357/5&53\\
D_9&763/10&8/5&751/10&61
\end{array}
\]
In the right-hand table, $d=\dim G$ and $q=d-R-\delta$ except in the
$B_2,C_2$ row, where the cap threshold is $T=9$ and $q=T-\delta$.
\end{proposition}

\begin{proof}
For a polynomial row, let
\[
 A=m_{2k+1}-Tm_{2k},\qquad
 B=m_{4k+2}-2Tm_{4k+1}+T^2m_{4k}.
\]
Lemma~\ref{lem:polynomial-one-sided-tail}, with $P(x)=x^k$, gives
$\mathbf P\{X>T\}\geq A^2/B$ whenever $A,B>0$.  The exact verifier cited
below regenerates all these moments from
Lemmas~\ref{lem:bounded-littlewood-classical-moments} and
\ref{lem:two-power-coefficient-functional} and checks, in reduced GMP
rational arithmetic,
\[
 A>0,\qquad B>0,\qquad
 \frac{A^2}{B}(1-\delta^{-2})(T-\delta)^K>(2C_G)^K.
\]
Lemma~\ref{lem:one-sided-cap-full-q3} proves the six polynomial rows.

For $B_2$ and $C_2$, use the type-$B_2$ positive roots
$\theta_1,\theta_2,\theta_1+\theta_2,\theta_1-\theta_2$.  On
$[0,1/4]^2$, the inequality $\cos u\geq1-u^2/2$ gives $X>9$.
Moreover, $4\sin^2(u/2)\geq u^2/4$ and $\pi<4$ give
\[
\begin{aligned}
 \mathbf P\{X>9\}
 &\geq\frac1{512}\left(\frac14\right)^4
 \int_0^{1/4}\!\int_0^{1/4}
 \theta_1^2\theta_2^2(\theta_1+\theta_2)^2
 (\theta_1-\theta_2)^2\,d\theta_1d\theta_2\\
 &=\frac1{512}\left(\frac14\right)^{14}
 \left(\frac1{21}-\frac2{25}+\frac1{21}\right)
 =\frac1{9019431321600}.
\end{aligned}
\]
The two rows have the same adjoint-character law because their Lie algebras
are isomorphic and Lemma~\ref{lem:full-central-cover} removes the global-form
choice.

For the remaining four cap rows, apply
Lemma~\ref{lem:full-q3-classical-weyl-cap}.  Its exact rational lower bound
$\widehat p_G$ is obtained as follows:
\[
\begin{aligned}
 \widehat p_{B_6}
 &=\frac{(1-R/24)^2\prod_{i=1}^6(2i)!}
 {2^6 6!\,2^6(44/7)^3 39!}
 \left(\frac R{22}\right)^{39},\\
 \widehat p_{D_6}
 &=\frac{(1-R/24)^2(2!4!6!8!10!6!)}
 {2^5 6!\,(44/7)^3 33!}
 \left(\frac R{20}\right)^{33},\\
 \widehat p_{D_7}
 &=\frac{(1-R/24)^2(2!4!6!8!10!12!7!)}{2^6 7!}
 \frac{4^{46}46!}{12(22/7)^4 92!}
 \left(\frac R{24}\right)^{45}\frac{433}{500},\\
 \widehat p_{D_9}
 &=\frac{(1-R/192)^{16}(2!4!6!8!10!12!14!16!9!)}{2^8 9!}
 \frac{(7/5)4^{77}77!}{(44/7)^5 154!}
 \left(\frac R{32}\right)^{76}\frac{193}{125}.
\end{aligned}
\]
Here $R$ has the row-specific value in the table.  These are strict lower
bounds in~\eqref{eq:full-q3-classical-cap-mass}: besides $\pi<22/7$, the
odd-dimensional rows use
\[
 \sqrt2<\frac32,\quad \sqrt{\frac34}>\frac{433}{500}
 \qquad(D_7),
\]
and
\[
 \sqrt2<\frac{10}{7},\quad
 \sqrt{\frac{763}{320}}>\frac{193}{125}
 \qquad(D_9),
\]
all verified by squaring positive rationals.  The sine factors are
\eqref{eq:full-q3-product-sine-small} in the first three rows and
\eqref{eq:full-q3-product-sine-long} in $D_9$.
The same exact ledger checks $R<9\kappa$ in all four rows; since
$\pi>3$, this places the Euclidean ball strictly inside the standard
eigenangle chart.

Finally, the same GMP verifier checks for each cap row
\[
 \widehat p_G(1-\delta^{-2})q^K>(2C_G)^K,
\]
using $C_{B_2}=C_{C_2}=2$, $C_{B_6}=C_{D_6}=6$, $C_{D_7}=5$, and
$C_{D_9}=7$.  Lemma~\ref{lem:one-sided-cap-full-q3} completes the proof.
Every rational comparison in this proof is implemented in
\begin{center}
\path{character_ring_iter/verify_full_q3_bcd_bounded_littlewood_gmp.cpp}.
\end{center}
The accepted transcript
\path{certificates/full_q3/fullq3bcd0010_machine_c.log} asserts the exact
scope of six polynomial rows and six rational-cap rows, prints a positive
reduced rational margin for each, and ends with exit status zero and the
terminal marker
\texttt{FULL\_Q3\_BCD\_BOUNDED\_LITTLEWOOD VERIFICATION: ALL PASS}.
\end{proof}

\begin{proposition}[Directed exact-MGF tails in the residual rows]
\label{prop:bcd-remaining-directed-tails}
For the ranks in the following displays, let $K_b$ be the corresponding
entry:
\[
\begin{array}{c|rrrrrrrrrrr}
b&7&8&9&10&11&12&13&14&15&16&17\\ \hline
K_b(B)&53&51&53&57&57&61&61&45&45&45&45
\end{array}
\]
\[
\begin{array}{c|rrrrrrrrrrrrr}
b&4&5&6&7&8&9&10&11&12&13&14&15&16\\ \hline
K_b(C)&61&55&39&39&49&49&53&55&53&57&59&55&55
\end{array}
\]
\[
\begin{array}{c|rrrrrrrrrrrr}
b&17&18&19&20&21&22&23&24&25&26&27&28\\ \hline
K_b(C)&57&45&45&45&47&47&49&51&53&55&57&59
\end{array}
\]
and
\[
\begin{array}{c|rrrrrrrrrrrr}
b&8&10&11&12&13&14&15&16&17&18&19&20\\ \hline
K_b(D)&57&51&57&45&45&45&45&43&43&43&45&45
\end{array}
\]
\[
\begin{array}{c|rrrrrrrrrr}
b&21&22&23&24&25&26&27&28&29&30\\ \hline
K_b(D)&45&45&45&47&45&47&47&49&47&49.
\end{array}
\]
Then, in every displayed row,
\[
 I^G_{a,n}>0
 \qquad(a\geq1,\ n\text{ odd},\ 2a+n\geq K_b).
\]
\end{proposition}

\begin{proof}
For each row, the certificate chooses exact rational numbers
$r,q,\delta>0$, sets $c=rb$ and
\[
 t=\sqrt{2c+2\delta+q},
\]
and applies Lemma~\ref{lem:finite-layered-mgf-recovery} with eight layers to
the exact determinants in
Lemma~\ref{lem:exact-defining-trace-determinants}.  In type $B$ the positive
and negative defining-trace tails are evaluated separately; in types $C$
and $D$ symmetry doubles the positive tail.  Let $V$ be the resulting lower
bound for $\mathbf P\{|T_1|\geq t\}$.

Rains' power-two decompositions
\cite[Theorem~1.5, equations~(22), (23), and (36)]{RainsPowerMaps} give an
upper bound $U$ for $\mathbf P\{\tau_G>q\}$ by exponential Markov.  In type
$B$ the bound is
\[
 U\leq\exp\{-\lambda(q-1)+\lambda^2\}.
\]
In types $C,D$, the MGF of $\tau_G-1$ is a product of two exact
orthogonal-component MGFs from
Lemma~\ref{lem:exact-defining-trace-determinants}, with the component sizes
and signs specified by the cited decomposition.  Thus the event $A$ in
Lemma~\ref{lem:all-minus-trace-tail} has probability at least $V-U$.
Every row certificate checks $V-U>0$.

Chebyshev gives $\mathbf P\{|X|<\delta\}\geq1-\delta^{-2}$.  At the
displayed onset $K=K_b$, the certificate bounds
\[
 \mathbf E(\tau_G)_+^K
 \leq\min\left\{d_T^K,
 \left(\frac{K}{ev}\right)^K\mathbf Ee^{v\tau_G}\right\},
 \qquad
 v=\frac{\sqrt{1+8K}-1}{4},
\]
where $d_T=2b+1$ in type $B$ and $d_T=2b$ in types $C,D$; the MGF is the
same source-audited product just described.  It then checks
\begin{equation}
\label{eq:remaining-directed-tail-test}
 (1-\delta^{-2})(V-U)c^K>\mathbf E(\tau_G)_+^K,
 \qquad c>2C_G.
\end{equation}
Lemma~\ref{lem:all-minus-trace-tail} propagates each row from $K$ to every
later odd total degree.

The proof companion is
\begin{center}
\path{character_ring_iter/verify_full_q3_bcd_low_tail_mpfr.cpp}.
\end{center}
It evaluates each Bessel series
\[
 I_k(2s)=\sum_{m\geq0}\frac{s^{2m+k}}{m!(m+k)!}
\]
through $180$ recurrence updates and encloses the remaining positive tail
by a geometric series.  It evaluates every determinant by interval Cholesky;
the matrices are Gram matrices for positive Jacobi weights, and the program
rejects a nonpositive pivot.  Every rational operation, square root,
exponential, logarithm, Bessel bound, and Cholesky operation is rounded
outward at $384$-bit MPFR precision.  The raw certificate prints the full
rational schedule for each row before its result and cross-checks all onsets
against the independent header
\path{character_ring_iter/full_q3_bcd_remaining_data.hpp}.
It verifies all $58$ rows; its least directed logarithmic margin in
\eqref{eq:remaining-directed-tail-test} is
\[
 0.0971194051797054477804174\ldots>0
 \quad\text{at }C_{15}.
\]
Its terminal marker is
\texttt{FULL\_Q3\_BCD\_LOW\_TAIL VERIFICATION: ALL PASS}.
The complete transcript is
\begin{center}
\path{certificates/full_q3/fullq3bcdlowtail0001_optimus.log},
\end{center}
authenticated by its adjacent SHA-256 file.
\end{proof}

\begin{proposition}[Exact bounded-Littlewood residual certificate]
\label{prop:bcd-bounded-littlewood-residual-certificate}
For a row $G=B_b,C_b,D_b$, put
\[
 s_G=\begin{cases}2b+1,&G=B_b,C_b,\\ b-1,&G=D_b.\end{cases}
\]
In the rows
\[
 B_2,\ldots,B_{17},\qquad
 C_2,\ldots,C_{28},\qquad
 D_4,\ldots,D_{30},
\]
let $K_G$ be the tail onset assigned to that row in
Proposition~\ref{prop:bcd-remaining-exact-tails} or
Proposition~\ref{prop:bcd-remaining-directed-tails}.  Then
\[
 I^G_{a,n}>0
 \quad\text{whenever}\quad
 a\geq2,\quad n\text{ is odd},\quad s_G<2a+n<K_G.
\]
This domain contains exactly $12{,}993$ pairs.  Its smallest value is
\[
 I^{D_5}_{2,1}=50.
\]
\end{proposition}

\begin{proof}
We describe the exact finite calculation and why it evaluates the formulas
in the preceding lemmas.  After applying $\pi_2$, substitute
\[
 p=t,\qquad q=zt^2.
\]
For the homogeneous degree-$2j$ part of any of
$\mathcal B_b,\mathcal C_b,\mathcal D_b$, let $P_{G,j}(z)$ be the
coefficient of $t^{2j}$.  It has degree at most $j$, and
Lemma~\ref{lem:two-power-coefficient-functional} becomes
\begin{equation}
\label{eq:bounded-littlewood-interpolation-functional}
 m_j^G=j!\sum_{a=0}^{j}(2(j-a)-1)!![z^a]P_{G,j}(z).
\end{equation}
Thus the values at any $j+1$ distinct rational points determine the moment
exactly.

The verifier
\begin{center}
\path{character_ring_iter/verify_full_q3_bcd_bounded_littlewood_gmp.cpp}
\end{center}
works simultaneously through $j=59$.  For each of the integer nodes
$z=0,1,\ldots,59$, it uses the recurrences in
Lemma~\ref{lem:two-power-coefficient-functional} to form the specialized
$e_r,h_r,f_r^e,f_r^h$ as truncated power series through degree $118$, then
evaluates every leading determinant in
Lemma~\ref{lem:bounded-littlewood-classical-moments}.  Lagrange interpolation
applies the functional in
\eqref{eq:bounded-littlewood-interpolation-functional}.  All of this is done
in exact prime fields; the primes exceed $2^{30}$, and therefore every
integer divided by the recurrences and every nonzero node difference is
invertible.

The $20$ prime-field evaluations are independent OpenMP tasks.  Their
residues are merged in one fixed descending-prime order by exact GMP CRT,
giving a modulus of $620$ bits.  If $d_G=\dim G$, invariant multiplicity and
$|\chiad|\leq d_G$ give the elementary reconstruction bound
\[
 0\leq m_j^G\leq d_G^j.
\]
The program chooses primes until their product is strictly larger than the
largest such bound it consumes, and rejects any reconstructed moment outside
its own bound.  Hence the CRT representative is the unique integer moment,
not a probable reconstruction.

The program next regenerates the stable recurrence, checks every stable
prefix, checks the $B_2=C_2$ Lie-isomorphism row, and substitutes the moments
into the exact GMP sum of Lemma~\ref{lem:full-moment-form}.  The row ledger is
\path{character_ring_iter/full_q3_bcd_remaining_data.hpp}; it asserts the
stable endpoint, tail onset, and moment endpoint for every one of the $70$
rows in the statement.  The verifier rejects a nonpositive hierarchy value,
an incorrect row ledger, or a scope other than $12{,}993$.

The accepted machine-C transcript is
\begin{center}
\path{certificates/full_q3/fullq3bcd0010_machine_c.log}.
\end{center}
It authenticates the verifier and row-ledger sources, reports all
$12{,}993$ values positive with the stated minimum, ends with
\texttt{FULL\_Q3\_BCD\_BOUNDED\_LITTLEWOOD VERIFICATION: ALL PASS}, and
records wrapper exit status zero.  Its adjacent SHA-256 file authenticates
the raw transcript.  As an implementation-independent control, the earlier
reverse-Pieri verifier agrees exactly on all $2{,}845$ moments and all
$7{,}484$ residual pairs through degree $40$; the accepted proof, however,
uses the determinant identities above and the complete degree-$59$ replay.
\end{proof}

\begin{proposition}[Exact $B_{18}$--$B_{21}$ and
$D_{31}$--$D_{70}$ residual certificate]
\label{prop:bd-finite-residual-certificate}
For every $B_b$ with $18\leq b\leq21$ and every $D_b$ with
$31\leq b\leq70$,
\[
 I^G_{a,n}>0\qquad(a\geq2,\ n\text{ odd}).
\]
Thus the complete generator hierarchy holds in all $44$ rows.
\end{proposition}

\begin{proof}
Theorem~\ref{thm:stable-bcd-full} supplies the stable prefix and
Proposition~\ref{prop:bcd-low-tail} supplies the tail.  It remains to check
the odd total degrees strictly between them.  The case $a=1$ is
Theorem~\ref{thm:two-minus-import}, so only $a\geq2$ remains.  Lemma
\ref{lem:finite-orthogonal-adjoint-moments} supplies every moment in that
finite box.  In type $D$, every consumed degree is at most $2b$, so no
unlisted even-partition overflow occurs.

The exact C++/GMP verifier
\begin{center}
\path{character_ring_iter/verify_full_q3_bd_residual_gmp.cpp}
\end{center}
regenerates $s_0,\ldots,s_{71}$ from
\[
 s_{r+1}=r s_r+r s_{r-1}-\binom r2s_{r-2},
\]
and rejects any disagreement with the vendored row.  For each required
moment it enumerates precisely the partitions in Lemma
\ref{lem:finite-orthogonal-adjoint-moments}; a reverse
vertical-two-strip Pieri traversal computes all $c_j(\lambda)$ with GMP
integers.  The traversals for distinct moment degrees are independent and
are distributed over OpenMP workers; their arithmetic and output are
unchanged by the worker count.  The program then substitutes the regenerated
moments in Lemma~\ref{lem:full-moment-form} for every residual pair and
rejects a nonpositive value or an incorrect scope count.

There are exactly $137$ type-$B$ and $4732$ type-$D$ residual pairs.  The
machine-C replay checks all $4869$ and reports the minimum
\[
 I^{D_{31}}_{8,15}
 =1272988116891284109857142182380802>0.
\]
The raw transcript is
\begin{center}
\path{certificates/full_q3/fullq3bcd0002_machine_c.log},
\end{center}
authenticated by its adjacent SHA-256 file.  Its terminal marker is
\texttt{FULL\_Q3\_B18\_D31\_RESIDUAL VERIFICATION: ALL PASS}.
This proves every residual inequality and hence the proposition.
\end{proof}

\begin{corollary}[Completed middle orthogonal rows]
\label{cor:completed-middle-orthogonal-rows}
The complete generator hierarchy holds for
\[
 B_b\quad(b\geq18),\qquad D_b\quad(b\geq31).
\]
\end{corollary}

\begin{proof}
Combine Propositions~\ref{prop:bcd-high-rank-full-tail},
\ref{prop:bcd-middle-rank-full-tail},
\ref{prop:bcd-low-tail}, and
\ref{prop:bd-finite-residual-certificate} with
Theorem~\ref{thm:stable-bcd-full}.
\end{proof}

\begin{theorem}[Full hierarchy in types $B$, $C$, and $D$]
\label{thm:bcd-full-hierarchy}
Let $G$ be a compact connected group with simple Lie algebra of type
$B_b$ $(b\geq2)$, $C_b$ $(b\geq2)$, or $D_b$ $(b\geq4)$.  Then
\[
 I^G_{a,n}\geq0\qquad(a,n\geq0).
\]
Consequently Ginibre's full condition $Q_3$ holds on $\Qcone_G$.
\end{theorem}

\begin{proof}
First consider the low rows $B_2,\ldots,B_{17}$,
$C_2,\ldots,C_{28}$, and $D_4,\ldots,D_{30}$.  The all-plus and even-$n$
sectors are automatic by Theorem~\ref{thm:full-cone-reduction}, while
Theorem~\ref{thm:two-minus-import} supplies $a=1$.  Let $a\geq2$, let $n$
be odd, and put $L=2a+n$.  Theorem~\ref{thm:stable-bcd-full} applies when
$L\leq s_G$.  Proposition
\ref{prop:bcd-bounded-littlewood-residual-certificate} applies when
$s_G<L<K_G$.  At $L\geq K_G$, the row is covered by either
Proposition~\ref{prop:bcd-remaining-exact-tails} or
Proposition~\ref{prop:bcd-remaining-directed-tails}.  These three ranges are
exhaustive.

For $B_b$ with $b\geq18$ and $D_b$ with $b\geq31$, use
Corollary~\ref{cor:completed-middle-orthogonal-rows}.  For $C_b$ with
$29\leq b\leq295$, use
Proposition~\ref{prop:bcd-middle-rank-full-tail}; the same proposition covers
the still-needed middle ranges $22\leq b\leq295$ in type $B$ and
$71\leq b\leq297$ in type $D$.  All remaining higher ranks satisfy the
hypothesis of Proposition~\ref{prop:bcd-high-rank-full-tail}.  Thus the
generator hierarchy holds in every row.  Central-cover invariance gives every
compact connected global form, and
Theorem~\ref{thm:full-cone-reduction} promotes the generator to $\Qcone_G$.
\end{proof}

\begin{theorem}[Full adjoint-generated Ginibre $Q_3$]
\label{thm:full-adjoint-generated-q3}
Let $G$ be any compact connected Lie group with simple Lie algebra.  For every
integer $r\geq0$, every tuple $f_1,\ldots,f_r\in\Qcone_G$, and every choice
of signs $\varepsilon_i\in\{+1,-1\}$ containing an even number of minus
signs,
\[
 \int_{G\times G}\prod_{i=1}^{r}
 \bigl(f_i(g)+\varepsilon_i f_i(h)\bigr)
 \,d\mu_G(g)d\mu_G(h)\geq0.
\]
If the number of minus signs is odd, the integral is zero.  Hence
$\Qcone_G$ satisfies Ginibre's condition $Q_3$ with all of its original
finite-tuple and sign-pattern quantifiers.
\end{theorem}

\begin{proof}
Classification leaves type $A$, types $B,C,D$, and the five exceptional
types.  Theorem~\ref{thm:type-a-full-cone},
Theorem~\ref{thm:bcd-full-hierarchy}, and
Theorem~\ref{thm:exceptional-full-cone}, respectively, prove the full
generator hierarchy in those cases.  Lemma~\ref{lem:full-central-cover}
removes the choice of compact connected global form.  Finally,
Theorem~\ref{thm:full-cone-reduction} gives the asserted cone promotion and
the vanishing of every odd-minus sector.
\end{proof}

\begin{remark}[Exact sense in which the result matches Ginibre]
Theorem~\ref{thm:full-adjoint-generated-q3} matches the full quantifier
structure of Ginibre's 1970 condition $Q_3$ and uses his Proposition~1
multiplicative-cone promotion.  It does not replace $\Qcone_G$ by the cone of
all real continuous positive-definite functions.  Such a replacement is
impossible: Proposition~\ref{prop:all-pd-counterexample} gives the exact
$\mathrm{SO}(3)$ value $-8/279$.
\end{remark}

\section*{Code and data availability}

\subsection*{Certificate semantics and trusted base}

The computer-assisted results have three separate layers.  First, each
numbered proposition states the mathematical input formula, the finite
quantified domain, and the sign or interval comparison to be decided.  Second,
the named verifier reconstructs the inputs and evaluates that comparison; a
stored output value is not accepted as an input merely because it appears in a
transcript.  Third, the transcript records the source identities, scope
counts, arithmetic outputs, terminal marker, and process status of one
execution.  A transcript authenticates a completed run but is not, by itself,
the proof of the formula implemented by the verifier.

The trusted computational base is consequently explicit: a conforming C++
implementation, GMP integer and rational arithmetic, MPFR directed rounding,
and the short orchestration and manifest parsers named below.  No claim is
made that these programs are formally verified.  The exact stages avoid
probabilistic reconstruction: when prime fields and CRT are used, the proved
modulus bound makes the reconstructed integer unique.  The interval stages
accept a sign only when the outward-rounded lower endpoint is positive.
Independent recurrence checks, alternative low-degree implementations,
sanitizer runs, and mutation tests are redundancy controls; their stated
cutoffs are not represented as full second proofs of computations beyond
those cutoffs.

The audited computation package for Parts~I--III is rooted at the immutable
publication import commit
\begin{center}
\url{https://github.com/skypher/ginibre-q3-adjoint/tree/91363c3cebb37a44349f613ab0a5aa6dcd412af3/ginibre_q3}.
\end{center}
Its full identifier is
\begin{center}
\texttt{91363c3cebb37a44349f613ab0a5aa6dcd412af3}.
\end{center}
The immutable import commit controls the execution snapshot that produced the
expensive arithmetic transcripts.  The historical tag \texttt{v1.0.0}
contains that import binding, but it is not asserted to contain later
manuscript revisions.  For the present submission, the canonical source bytes
are the files in the submitted archive whose hashes are listed in
\path{full_q3_source_manifest.sha256}.  A public release of the final
submission must reproduce that manifest exactly.

The formal submission has three numbered parts and a formal detailed
supplement, supplied as three inseparable manuscript files:
\path{paper.pdf}, containing the compact Parts~I--II,
\path{paper_full.pdf}, containing their load-bearing correction-prefix
derivations, and \path{full_q3_extension.pdf}, containing the present
Part~III.  The archived
computation-bearing sources are \path{paper.tex}, \path{paper_full.tex}, and
\path{full_q3_extension.tex} in that commit.  The immutable import commit
authenticates the arithmetic snapshot.  The current publication sources are
bound directly by \path{full_q3_source_manifest.sha256} and are the inputs
rebuilt by the commands below.  The historical aggregate transcript predates
publication-only prose and metadata revisions; the fail-closed
\path{verify_full_q3_related_work_rebind.py} records that boundary and rejects
any change to an arithmetic-bearing source.  That rebind is a provenance test,
not a step required to construct the submitted source: the final files are
bound directly by the final-source manifest.  A final release archive must
preserve both the immutable execution snapshot and those exact final-source
bytes.
Part~III is not offered as a
standalone proof: its Theorem~\ref{thm:two-minus-import} imports Theorem~2.2
of Parts~I--II, which is included in the submission and archive.

The formal detailed supplement \path{paper_full.pdf} is rebuilt from the same
source.  Its numbered correction-prefix derivations are incoming proof nodes
for the residual $B/C$ closure; diagnostic and superseded material printed
there remains outside the dependency spine.  The compact contract,
accepted-certificate boundary, and family closures appear in
\path{paper.pdf}.

The exact replay contract, software requirements, host/thread guidance, and
terminal success markers are documented in \path{REPLAY.md} and
\path{FULL_Q3_DISTRIBUTED_REPLAY.md}.  The principal requirements are a
C++20 compiler with OpenMP, GMP, MPFR, Python~3, GNU Make, and a \LaTeX{}
installation.  From the repository root, the three aggregate commands are
\begin{center}
\texttt{python3 ginibre\_q3/run\_publication\_short\_audits.py}

\texttt{make -C ginibre\_q3 clean-room-replay}

\texttt{make -C ginibre\_q3 full-q3-extension}.
\end{center}
These commands implement three reproducibility tiers.  The first command runs
the document, dependency, interface, scope, moment-formula, cone-promotion,
$SO(3)$, and frontier-formula checks as one fail-closed short suite.  The
Part~III command regenerates its exact and
directed-rounding arithmetic; the archived two-worker run took 44 minutes 55
seconds.  The separate Parts~I--II clean-room command is the full
certificate-regeneration tier: its partitioned stages have two- or eight-hour
guards and require the machine-C resource envelope documented in
\path{REPLAY.md}.  Part~III does not silently substitute its faster audit for
that imported computation.

SHA-256 manifests authenticate every load-bearing source, accepted
transcript, generated ledger, and cited source artifact.  Compiler and
resource checks and author self-audits are integrity or redundancy tests; they
are not independent peer review and do not by themselves decide a sign.  The
mathematical certificate claims are the exact integer/rational or directed-rounding
comparisons described in the numbered results above.

Table~\ref{tab:certificate-map} gives the compact trusted-computation map.
Here ``exact'' means integer or rational arithmetic, including prime-field
evaluation followed by uniquely bounded GMP reconstruction; ``directed''
means outward-rounded 384-bit MPFR interval arithmetic.
\begin{table}[ht]
\centering
\scriptsize
\begin{tabular}{@{}L{.17\textwidth}L{.19\textwidth}L{.14\textwidth}L{.36\textwidth}@{}}
\hline
Result block & Program/transcript identifier & Arithmetic & Verified scope \\
\hline
Imported two-minus theorem
 & \texttt{distributed0001}
 & exact/directed
 & 200 replay stages; all compact connected simple types; every $n\geq0$ \\
Low and stable type $A$
 & \texttt{su2}, \texttt{su3}, \texttt{su45}, \texttt{sun-ge6}
 & exact GMP
 & all residual boxes and exact tail comparisons; 1,673 residual pairs \\
Exceptional types
 & \texttt{exceptional}
 & exact GMP
 & five tails and 1,265 residual pairs \\
High and middle classical tails
 & \texttt{bcdanalytic0003}
 & directed MPFR
 & all $C_G\geq296$ rows, 768 middle rows, and 44 supplemental rows \\
Remaining classical tails
 & \texttt{bcdlowtail0001}
 & directed MPFR
 & 58 finite-rank tail rows \\
Low classical residual boxes
 & \texttt{bcd0010}
 & exact prime/GMP
 & 70 rows and 12,993 residual inequalities \\
Finite $B_{18}$--$B_{21}$ and $D_{31}$--$D_{70}$ boxes
 & \texttt{bcd0002}
 & exact GMP
 & 4,869 residual inequalities \\
\hline
\end{tabular}
\caption{Map from the computer-assisted proof blocks to their accepted
program/transcript families and arithmetic boundaries.}
\label{tab:certificate-map}
\end{table}

The identifiers in the table resolve to the following authenticated
transcripts:
\begin{itemize}
\item \path{certificates/full_q3/distributed0001/fullq3distributed0001.log};
\item \path{certificates/full_q3/fullq3complete0002_optimus.log};
\item \path{certificates/full_q3/fullq3bcdanalytic0003_machine_c.log};
\item \path{certificates/full_q3/fullq3bcdlowtail0001_optimus.log};
\item \path{certificates/full_q3/fullq3bcd0010_machine_c.log}; and
\item \path{certificates/full_q3/fullq3bcd0002_machine_c.log}.
\end{itemize}

\begin{thebibliography}{99}

\bibitem{AdjointTwoMinus}
L.~P. Polzer,
The two-minus adjoint-character case of Ginibre's $Q_3$ condition,
Parts I--II: classification and exact certificates, journal manuscript with
formal detailed supplement, 2026,
\url{https://github.com/skypher/ginibre-q3-adjoint/tree/91363c3cebb37a44349f613ab0a5aa6dcd412af3/ginibre_q3}.

\bibitem{DiaconisShahshahani1994}
P.~Diaconis and M.~Shahshahani,
On the eigenvalues of random matrices,
\emph{J. Appl. Probab.} \textbf{31A} (1994), 49--62.

\bibitem{Ginibre1970}
J.~Ginibre,
General formulation of Griffiths' inequalities,
\emph{Comm. Math. Phys.} \textbf{16} (1970), 310--328.

\bibitem{Sylvester1980}
G.~S. Sylvester,
The Ginibre inequality,
\emph{Comm. Math. Phys.} \textbf{73} (1980), no.~2, 105--114.

\bibitem{Abdesselam2022}
A.~Abdesselam,
Non-Abelian correlation inequalities and stable determinantal polynomials,
arXiv:2207.07603, 2022.

\bibitem{ChevyrevGarban2025}
I.~Chevyrev and C.~Garban,
Villain action in lattice gauge theory,
\emph{J. Stat. Phys.} \textbf{192} (2025), Paper~38.

\bibitem{HKKO2025}
J.~Huh, J.~S. Kim, C.~Krattenthaler, and S.~Okada,
Bounded Littlewood identities with fixed number of odd rows or odd columns,
\emph{Electron. J. Combin.} \textbf{33} (2026), no.~3, Paper~P3.5.

\bibitem{Etingof2009}
P.~Etingof,
A uniform proof of the Macdonald--Mehta--Opdam identity for finite Coxeter
groups,
arXiv:0903.5084, 2009.

\bibitem{GaribaldiGuralnickRains2025}
S.~Garibaldi, R.~M.~Guralnick, and E.~M.~Rains,
Invariant derivations and trace bounds,
\emph{J. \'{E}c. polytech. Math.} \textbf{12} (2025), 1229--1287.

\bibitem{BPV2024}
J.~L. Bourjaily, M.~R. Plesser, and C.~Vergu,
The many colours of amplitudes,
arXiv:2412.21189.

\bibitem{Rains1997}
E.~M. Rains,
Increasing subsequences and the classical groups,
\emph{Electron. J. Combin.} \textbf{5} (1998), Research Paper 12.

\bibitem{RainsPowerMaps}
E.~M. Rains,
Images of eigenvalue distributions under power maps,
\emph{Probab. Theory Related Fields} \textbf{125} (2003), 522--538.

\bibitem{CJL2022}
C.~Courteaut, K.~Johansson, and G.~Lambert,
From Berry--Esseen to super-exponential,
arXiv:2204.03282.

\end{thebibliography}

\end{document}
